
\Data\ID{1003025102}
\Updated{2021-08-19 10:08:04}
\Level{A}
\SubArea{Statistics}
\LANGen{
\Question{
A car travels the first quarter of its journey at an average speed of \( 50\, \mathrm{kph} \), the second quarter at an average speed of \( 90\, \mathrm{kph} \), the third quarter at an average speed of \( 130\, \mathrm{kph} \) and the remaining one quarter at an average speed of \( 80\, \mathrm{kph} \). What is the average speed of the car during the journey? Round the result to two decimal places.}
{\( 77{.}97\, \mathrm{km}/\mathrm{h} \)}{\( 85{.}00\, \mathrm{km}/\mathrm{h} \)}{\( 87{.}50\, \mathrm{km}/\mathrm{h} \)}{\( 82{.}71\, \mathrm{km}/\mathrm{h} \)}{}{}}
\LANGpl{
\Question{
Droga samochodu została podzielna na cztery równe części. Samochód pierwszą część przejechał z prędkością  \( 50\, \mathrm{km}/\mathrm{h} \), drugą  część  przejechał z prędkością \( 90\, \mathrm{km}/\mathrm{h} \), trzecią część przejechał z prędkością \( 130\, \mathrm{km}/\mathrm{h} \), a czwartą z prędkością \( 80\, \mathrm{km}/\mathrm{h} \). Jaka jest średnia prędkość samochodu w czasie  podróży? Zaokrąglij wynik do dwóch miejsc po przecinku.}
{\( 77{,}97\, \mathrm{km}/\mathrm{h} \)}{\( 85{,}00\, \mathrm{km}/\mathrm{h} \)}{\( 87{,}50\, \mathrm{km}/\mathrm{h} \)}{\( 82{,}71\, \mathrm{km}/\mathrm{h} \)}{}{}}
\LANGcs{
\Question{
Auto jelo první čtvrtinu trasy průměrnou rychlostí \( 50\, \mathrm{km}/\mathrm{h} \), druhou čtvrtinu trasy průměrnou rychlostí \( 90\, \mathrm{km}/\mathrm{h} \), třetí čtvrtinu trasy průměrnou rychlostí \( 130\, \mathrm{km}/\mathrm{h} \) a čtvrtou čtvrtinu trasy průměrnou rychlostí \( 80\, \mathrm{km}/\mathrm{h} \). Jakou průměrnou rychlostí jelo auto? Výsledek zaokrouhlete na 2 desetinná místa.}
{\( 77{,}97\, \mathrm{km}/\mathrm{h} \)}{\( 85{,}00\, \mathrm{km}/\mathrm{h} \)}{\( 87{,}50\, \mathrm{km}/\mathrm{h} \)}{\( 82{,}71\, \mathrm{km}/\mathrm{h} \)}{}{}}
\LANGsk{
\Question{
Auto išlo prvú štvrtinu cesty priemernou rýchlosťou \( 50\, \mathrm{km}/\mathrm{h} \), druhú štvrtinu cesty priemernou rýchlosťou \( 90\, \mathrm{km}/\mathrm{h} \), tretiu štvrtinu cesty priemernou rýchlosťou \( 130\, \mathrm{km}/\mathrm{h} \) a štvrtú štvrtinu cesty priemernou rýchlosťou \( 80\, \mathrm{km}/\mathrm{h} \). Akou priemernou rýchlosťou išlo auto? Výsledok zaokrúhlite na 2 desatinné miesta.}
{\( 77{,}97\, \mathrm{km}/\mathrm{h} \)}{\( 85{,}00\, \mathrm{km}/\mathrm{h} \)}{\( 87{,}50\, \mathrm{km}/\mathrm{h} \)}{\( 82{,}71\, \mathrm{km}/\mathrm{h} \)}{}{}}
\LANGes{
\Question{
Un coche ha viajado un cuarto de distancia con una velocidad de \( 50\, \mathrm{km}/\mathrm{h} \), el segundo cuarto de la distancia con una velocidad de \( 90\, \mathrm{km}/\mathrm{h} \), el tercer cuarto de la distancia con la velocidad de \( 130\, \mathrm{km}/\mathrm{h} \) y el último cuarto con la velocidad de \( 80\, \mathrm{km}/\mathrm{h} \). ¿Cuál era la velocidad media del coche? Aproxima el resultado a 2 cifras decimales.}
{\( 77{,}97\, \mathrm{km}/\mathrm{h} \)}{\( 85{,}00\, \mathrm{km}/\mathrm{h} \)}{\( 87{,}50\, \mathrm{km}/\mathrm{h} \)}{\( 82{,}71\, \mathrm{km}/\mathrm{h} \)}{}{}}
\End

\Data\ID{1003025103}
\Updated{2021-02-19 12:02:49}
\Level{A}
\SubArea{Statistics}
\LANGen{
\Question{
Ten workers produce the same type of components. Two workers produce one component in \( 4 \) minutes, other three workers in \( 5 \) minutes, another one worker in \( 6 \) minutes, next  three workers in \( 7 \) minutes and the last one of them in \( 8 \) minutes. What is the average time needed to produce one component? Round the result to the nearest hundredth.}
{\( 5{.}49\,\mathrm{min} \)}{\( 5{.}50\, \mathrm{min} \)}{\( 5{.}65\, \mathrm{min} \)}{\( 5{.}80\, \mathrm{min} \)}{}{}}
\LANGpl{
\Question{
Dziesięciu pracowników wytwarza ten sam produkt. Dwóch pracowników robi to w \( 4 \) minuty, następnych trzech w \( 5 \) minut, kolejny pracownik w  \( 6 \) minut, następne trzy osoby w  \( 7 \) minut i ostatni w \( 8 \) minut. Jaki jest średni czas potrzebny do wyprodukowania jednego produktu? Zaokrąglij wynik do dwóch miejsc po przecinku.}
{\( 5{,}49\,\mathrm{min} \)}{\( 5{,}50\, \mathrm{min} \)}{\( 5{,}65\, \mathrm{min} \)}{\( 5{,}80\, \mathrm{min} \)}{}{}}
\LANGcs{
\Question{
Deset dělníků pracuje v dílně, v níž se vyrábí jeden typ součástky. Dva z nich dokážou jednu součástku vyrobit za \( 4 \) minuty, další tři k tomu potřebují \( 5 \) minut, jeden \( 6 \) minut, další tři \( 7 \) minut a poslední z nich \( 8 \) minut. Určete průměrný čas, za který dělníci zvládnou vyrobit tuto součástku. Výsledek zaokrouhlete na \( 2 \) desetinná místa.}
{\( 5{,}49\,\mathrm{min} \)}{\( 5{,}50\, \mathrm{min} \)}{\( 5{,}65\, \mathrm{min} \)}{\( 5{,}80\, \mathrm{min} \)}{}{}}
\LANGsk{
\Question{
Aký je priemerný čas, ktorý potrebuje jeden robotník na výrobu jednej súčiastky, ak z desiatich robotníkov dvaja z nich vyrobia jednu súčiastku už za \( 4 \) minúty, ďalší traja za \( 5 \) minút, jeden za \( 6 \) minút, ďalší traja za \( 7 \) minút a posledný z nich za \( 8 \) minút? Výsledok zaokrúhlite na \( 2 \) desatinné miesta.}
{\( 5{,}49\,\mathrm{min} \)}{\( 5{,}50\, \mathrm{min} \)}{\( 5{,}65\, \mathrm{min} \)}{\( 5{,}80\, \mathrm{min} \)}{}{}}
\LANGes{
\Question{
Diez obreros trabajan en un negocio produciendo un tipo de componente. Dos de ellos son capaces de producir un componente en \( 4 \) minutos, tres necesitan \( 5 \) minutos para producirlo, uno necesita \( 6 \) minutos, tres necesitan \( 7 \) minutos y el último \( 8 \) minutos. Averigua el tiempo medio que los obreros necesitan para producir esta componente. Aproxima el resultado a \( 2 \) cifras decimales.}
{\( 5{,}49\,\mathrm{min} \)}{\( 5{,}50\, \mathrm{min} \)}{\( 5{,}65\, \mathrm{min} \)}{\( 5{,}80\, \mathrm{min} \)}{}{}}
\End

\Data\ID{1003025104}
\Updated{2021-02-19 11:02:30}
\Level{A}
\SubArea{Statistics}
\LANGen{
\Question{
The annual production of a business is recorded in the following table. Find the compound annual growth rate (i.e. the average annual coefficient of the production growth, i.e. ratio that provides a constant growth rate over the time period) over the time period \( 2014 \) - \( 2017 \). Round the result to four decimal places.\[
\begin{array}{|c|c|c|c|c|} \hline \text{Year} & 2014 & 2015 & 2016 & 2017 \\\hline \text{Production (pcs)} & 20\: 000 & 20\: 400& 21\: 420 & 24\: 633 \\\hline
\end{array}\]}
{\( 1{.}0719 \)}{\( 1{.}0705 \)}{\( 1{.}0733 \)}{\( 1{.}0727 \)}{}{}}
\LANGpl{
\Question{
Oblicz  średni roczny współczynnik wzrostu produkcji w latach \( 2014 \) - \( 2017 \). Roczna produkcja jest podana w  tabeli.  Zaokrągli wynik do czterech miejsc po przecinku.\[
\begin{array}{|c|c|c|c|c|} \hline \text{Rok} & 2014 & 2015 & 2016 & 2017 \\\hline \text{Produkcja (szt)} & 20\: 000 & 20\: 400& 21\: 420 & 24\: 633 \\\hline
\end{array}\]}
{\( 1{,}0719 \)}{\( 1{,}0705 \)}{\( 1{,}0733 \)}{\( 1{,}0727 \)}{}{}}
\LANGcs{
\Question{
Určete průměrný roční koeficient růstu výroby v letech \( 2014 \) - \( 2017 \) v podniku, jehož roční výroba je zaznamenaná v tabulce. Výsledek zaokrouhlete na \( 4 \) desetinná místa.
\[
\begin{array}{|c|c|c|c|c|} \hline \text{Rok} & 2014 & 2015 & 2016 & 2017 \\\hline \text{Výroba (ks)} & 20\: 000 & 20\: 400& 21\: 420 & 24\: 633 \\\hline
\end{array}\]}
{\( 1{,}0719 \)}{\( 1{,}0705 \)}{\( 1{,}0733 \)}{\( 1{,}0727 \)}{}{}}
\LANGsk{
\Question{
Určte priemerný ročný koeficient rastu výroby v rokoch \( 2014 \) - \( 2017 \) v podniku, ktorého ročná výroba je zaznamenaná v tabuľke. Výsledok zaokrúhlite na \( 4 \) desatinné miesta.
\[
\begin{array}{|c|c|c|c|c|} \hline \text{Rok} & 2014 & 2015 & 2016 & 2017 \\\hline \text{Výroba (ks)} & 20\: 000 & 20\: 400& 21\: 420 & 24\: 633 \\\hline
\end{array}\]}
{\( 1{,}0719 \)}{\( 1{,}0705 \)}{\( 1{,}0733 \)}{\( 1{,}0727 \)}{}{}}
\LANGes{
\Question{
Calcula el coeficiente de creimiento anual de la producción en los años  \( 2014 \) - \( 2017 \) en un negocio, cuya producción está en la tabla. Aproxima el resultado a \( 4 \) cifras decimales.
\[
\begin{array}{|c|c|c|c|c|} \hline \text{Año} & 2014 & 2015 & 2016 & 2017 \\\hline \text{Producción (piezas)} & 20\: 000 & 20\: 400& 21\: 420 & 24\: 633 \\\hline
\end{array}\]}
{\( 1{,}0719 \)}{\( 1{,}0705 \)}{\( 1{,}0733 \)}{\( 1{,}0727 \)}{}{}}
\End

\Data\ID{1003029501}
\Updated{2021-02-19 12:02:34}
\Level{A}
\SubArea{Statistics}
\LANGen{
\Question{
In a toy-factory four workers make the same toys manually. In one \( 8 \) hour shift the first worker produced \( 12 \) toys, the second \( 10 \) toys, the third \( 16 \) toys and the fourth \( 12 \) toys. What was the average time to produce a toy on that work shift?}
{\( 38\,\mathrm{min}\ 24\,\mathrm{s} \)}{\( 38\,\mathrm{min}\ 40\,\mathrm{s} \)}{\( 39\,\mathrm{min}\ 30\,\mathrm{s} \)}{\( 38\,\mathrm{min}\ 58\,\mathrm{s} \)}{}{}}
\LANGpl{
\Question{
Czterech pracowników fabryki zabawek ręcznie wykonuje zabawki. Podczas jednej  \( 8 \) godzinnej zmiany pierwszy pracownik wykonał \( 12 \) zabawek, drugi \( 10 \) zabawek, trzeci \( 16 \) zabawek a czwarty \( 12 \) zabawek. Jaki był  średni czas wykonania zabawki w czasie tej zmiany?}
{\( 38\,\mathrm{min}\ 24\,\mathrm{s} \)}{\( 38\,\mathrm{min}\ 40\,\mathrm{s} \)}{\( 39\,\mathrm{min}\ 30\,\mathrm{s} \)}{\( 38\,\mathrm{min}\ 58\,\mathrm{s} \)}{}{}}
\LANGcs{
\Question{
Čtyři dělníci v dílně ručně vyrábí stejné hračky. Během \( 8 \) hodinové pracovní doby vyrobil první dělník  \( 12 \) hraček, druhý \( 10 \) hraček, třetí \( 16 \) hraček a čtvrtý \( 12 \) hraček. Jaká byla v tento den průměrná doba potřebná k výrobě jedné hračky?}
{\( 38\,\mathrm{min}\ 24\,\mathrm{s} \)}{\( 38\,\mathrm{min}\ 40\,\mathrm{s} \)}{\( 39\,\mathrm{min}\ 30\,\mathrm{s} \)}{\( 38\,\mathrm{min}\ 58\,\mathrm{s} \)}{}{}}
\LANGsk{
\Question{
V továrni štyria robotníci ručne vyrábajú rovnaké hračky. V priebehu celej \( 8 \)-hodinovej pracovnej doby vyrobili: prvý  \( 12 \) hračiek, druhý \( 10 \) hračiek, tretí \( 16 \) hračiek a štvrtý tiež \( 12 \) hračiek. Aká je priemerná doba výroby jednej hračky?}
{\( 38\,\mathrm{min}\ 24\,\mathrm{s} \)}{\( 38\,\mathrm{min}\ 40\,\mathrm{s} \)}{\( 39\,\mathrm{min}\ 30\,\mathrm{s} \)}{\( 38\,\mathrm{min}\ 58\,\mathrm{s} \)}{}{}}
\LANGes{
\Question{
Cuatro obreros producen juguetes en un negocio. Durante un turno de trabajo de \( 8 \) horas el primer obrero produció \( 12 \) juguetes, el segundo \( 10 \) juguetes, el tercero \( 16 \) juguetes y el cuarto \( 12 \) juguetes. ¿Cuál era el tiempo medio necesitado para producir un juguete ese día?}
{\( 38\,\mathrm{min}\ 24\,\mathrm{s} \)}{\( 38\,\mathrm{min}\ 40\,\mathrm{s} \)}{\( 39\,\mathrm{min}\ 30\,\mathrm{s} \)}{\( 38\,\mathrm{min}\ 58\,\mathrm{s} \)}{}{}}
\End

\Data\ID{1003029502}
\Updated{2021-02-19 12:02:16}
\Level{A}
\SubArea{Statistics}
\LANGen{
\Question{
Two employees work in a factory workroom. The first one completed an assigned task in \( 20 \) minutes, the second one completed the same task in \( 10 \) minutes. We are interested in the average time of the task completion. What type of the average do we need to use?}
{Arithmetic mean}{Harmonic mean}{Geometric mean}{Weighted arithmetic mean}{}{}}
\LANGpl{
\Question{
Dwóch pracowników pracuje w  fabryce. Pierwszy wykonał przydzielone zadanie w ciągu  \( 20 \) minut, drugi wykonał to samo zadanie w ciągu  \( 10 \) minut. Interesuje nas średni czas wykonania zadania. Jakiego rodzaju średniej potrzebujemy?}
{Średnia arytmetyczna}{Średnia harmoniczna}{Średnia geometryczna}{Ważona średnia harmoniczna}{}{}}
\LANGcs{
\Question{
V dílně pracují dva zaměstnanci. První splní danou úlohu za \( 20 \) minut, druhý splní tu samou úlohu za \( 10 \) minut. Zajímá nás, kolik minut průměrně trvá splnění dané úlohy. Jaký typ průměru se pro výpočet použije?}
{Aritmetický průměr}{Harmonický průměr}{Geometrický průměr}{Vážený aritmetický průměr}{}{}}
\LANGsk{
\Question{
V dielni pracujú dvaja zamestnanci. Prvý splní danú úlohu za \( 20 \) minút, druhý splní tú istú úlohu za \( 10 \) minút. Zaujíma nás, koľko priemerne minúť trvá splnenie danej úlohy. Aký typ priemeru sa pre výpočet použije?}
{Aritmetický priemer}{Harmonický priemer}{Geometrický priemer}{Vážený aritmetický priemer}{}{}}
\LANGes{
\Question{
En un negocio trabajan dos obreros. El primero cumple su tarea en \( 20 \) minutos, el segundo la hace en \( 10 \) minutos. Nos interesa cuántos minutos son necesarios para hacer la tarea en promedio. 
¿Qué tipo de media vamos a calcular?}
{Media aritmética}{Media armónica}{Media geométrica}{Media aritmética ponderada}{}{}}
\End

\Data\ID{1003029503}
\Updated{2021-02-19 12:02:43}
\Level{A}
\SubArea{Statistics}
\LANGen{
\Question{
Paul drove the first halfway of a test track at a constant speed of \( 20\,\mathrm{kph} \) and the second halfway at a constant speed of \( 30\,\mathrm{kph} \). We wish to know Paul's average speed. What type of the average do we need to use?}
{Harmonic mean}{Arithmetic mean}{Geometric mean}{Weighted arithmetic mean}{}{}}
\LANGpl{
\Question{
Paweł przejechał pierwszą połowę toru testowego ze stałą  prędkością   \( 20\,\mathrm{km/h} \), drugą połowę toru ze stałą prędkością \( 30\,\mathrm{km/h} \). Chcemy poznać średnią prędkość uzyskaną przez  Pawła. Jakiego rodzaju średnią należy użyć do obliczenia tej prędkości?}
{Średnia harmoniczna}{Średnia arytmetyczna}{Średnia geometryczna}{Ważona średnia arytmetyczna}{}{}}
\LANGcs{
\Question{
Pavel jel při testování auta na zkušební dráze první polovinu cesty stálou rychlostí \( 20\,\mathrm{km/h} \) a druhou polovinu cesty rychlostí \( 30\,\mathrm{km/h} \). Zajímá nás průměrná Pavlova rychlost. Jaký typ průměru se pro výpočet použije?}
{Harmonický průměr}{Aritmetický průměr}{Geometrický průměr}{Vážený aritmetický průměr}{}{}}
\LANGsk{
\Question{
Pavol jazdil na testovacej dráhe polovicu cesty stálou rýchlosťou \( 20\,\mathrm{km/h} \) a druhú polovicu cesty rýchlosťou \( 30\,\mathrm{km/h} \). Zaujíma nás priemerná Pavlova rýchlosť. Aký typ priemeru sa pre výpočet použije?}
{Harmonický priemer}{Aritmetický priemer}{Geometrický priemer}{Vážený aritmetický priemer}{}{}}
\LANGes{
\Question{
Durande la prueba de un coche Pablo ha ido con la velocidad \( 20\,\mathrm{km/h} \) el primer medio de su ruta y el otro medio ha ido con la velocidad \( 30\,\mathrm{km/h} \). Estamos interesados en su velocidad media. ¿Qué tipo de media vamos a calcular?}
{Media armónica}{Media aritmética}{Media geométrica}{Media aritmética ponderada}{}{}}
\End

\Data\ID{1003029504}
\Updated{2021-02-19 12:02:27}
\Level{A}
\SubArea{Statistics}
\LANGen{
\Question{
In five consecutive years the annual production growth was \( 1\% \), \( 8\% \), \( 0\% \), \( 4\% \) and \( 1\% \) respectively. Find the compound annual growth rate of this five-year period. Round the result to two decimal places.}
{\( 2{.}76\% \)}{\( 2{.}75\% \)}{\( 2{.}72\% \)}{\( 2{.}80\% \)}{}{}}
\LANGpl{
\Question{
W pięciu kolejnych latach roczny wzrost produkcji wyniósł odpowiednio  \( 1\% \), \( 8\% \), \( 0\% \), \( 4\% \) and \( 1\% \). Wskaż złożoną roczną stopę wzrostu w tym pięcioletnim okresie. Zaokrąglij wynik do dwóch miejsc po przecinku.}
{\( 2{,}76\% \)}{\( 2{,}75\% \)}{\( 2{,}72\% \)}{\( 2{,}80\% \)}{}{}}
\LANGcs{
\Question{
V pěti po sobě jdoucích letech byl roční růst výroby \( 1\:\% \), \( 8\:\% \), \( 0\:\% \), \( 4\:\% \) a \( 1\:\% \). Určete průměrný procentuální růst výroby za toto pětileté období. Výsledek zaokrouhlete na dvě desetinná místa.}
{\( 2{,}76\:\% \)}{\( 2{,}75\:\% \)}{\( 2{,}72\:\% \)}{\( 2{,}80\:\% \)}{}{}}
\LANGsk{
\Question{
V piatich po sebe idúcich rokoch bol ročný rast výroby \( 1\% \), \( 8\% \), \( 0\% \), \( 4\% \) and \( 1\% \). Určite priemerný percentuálny rast výroby za toto päťročné obdobie. Výsledok zaokrúhlite na dve desatinné miesta.}
{\( 2{,}76\% \)}{\( 2{,}75\% \)}{\( 2{,}72\% \)}{\( 2{,}80\% \)}{}{}}
\LANGes{
\Question{
En cinco años consecutivos el crecimiento de la producción era \( 1\:\% \), \( 8\:\% \), \( 0\:\% \), \( 4\:\% \) y \( 1\:\% \) respectivamente. Averigua el crecimiento anual medio durante este cinco años. Aproxima el resultado a dos cifras decimales.}
{\( 2{,}76\:\% \)}{\( 2{,}75\:\% \)}{\( 2{,}72\:\% \)}{\( 2{,}80\:\% \)}{}{}}
\End

\Data\ID{1003029505}
\Updated{2021-02-19 01:02:13}
\Level{A}
\SubArea{Statistics}
\LANGen{
\Question{
Your stock portfolio had the following returns in the five consecutive years: \( 10\% \), \( -20\% \), \( 0\% \), \( 10\% \), \( 20\% \) (the minus sign denotes the financial loss). We are interested in five years compound annual growth rate of the portfolio. What type of average do we need  to use?}
{Geometric mean}{Arithmetic mean}{Harmonic mean}{Weighted arithmetic mean}{}{}}
\LANGpl{
\Question{
Twoje portfel akcji miał następujące zwroty w ciągu pięciu następujących po sobie lat: \( 10\% \), \( -20\% \), \( 0\% \), \( 10\% \), \( 20\% \)  (znak minus oznacza stratę finansową). Jesteśmy zainteresowani pięcioletnią złożoną roczną stopą wzrostu portfela. Jakiego rodzaju średniej potrzebujemy?}
{Średnia geometryczna}{Średnia arytmetyczna}{Średnia harmoniczna}{Ważona średnia arytmetyczna}{}{}}
\LANGcs{
\Question{
Vaše portfólio akcií vykazovalo v pěti po sobě následujících letech roční výnosy: \( 10\:\% \), \( -20\:\% \), \( 0\:\% \), \( 10\:\% \), \( 20\:\% \) (záporné znaménko znamená ztrátu). Zajímá nás průměrný procentuální roční výnos za sledované období. Jaký typ průměru se pro výpočet použije?}
{Geometrický průměr}{Aritmetický průměr}{Harmonický průměr}{Vážený aritmetický průměr}{}{}}
\LANGsk{
\Question{
Vaše portfólio akcií vykazovalo v piatich po sebe nasledujúcich rokoch takéto ročné výnosy: \( 10\% \), \( -20\% \), \( 0\% \), \( 10\% \), \( 20\% \) (záporné znamienko znamená stratu). Zaujíma nás priemerný percentuálny ročný výnos za sledované obdobie. Aký typ priemeru sa pre výpočet použije?}
{Geometrický priemer}{Aritmetický priemer}{Harmonický priemer}{Vážený aritmetický priemer}{}{}}
\LANGes{
\Question{
Tu cartera de acciones mostró ganancias de\( 10\:\% \), \( -20\:\% \), \( 0\:\% \), \( 10\:\% \), \( 20\:\% \) en cinco años consecutivos (el signo menos significa una perdida). Estamos interesados en la ganancia anual media en porcentaje. ¿Qué tipo de media vamos a calcular?}
{Media geométrica}{Media aritmética}{Media armónica}{Media aritmética ponderada}{}{}}
\End

\Data\ID{1003029506}
\Updated{2021-02-19 12:02:29}
\Level{A}
\SubArea{Statistics}
\LANGen{
\Question{
Let us assume that the value of an unique artwork will increase \( 1{.}5 \) times in the first year, \( 1{.}2 \) times in the second year and \( 1{.}9 \) times in the third year. We are interested in the compound annual growth rate in these three years. What type of average do we need to use?}
{Geometric mean}{Arithmetic mean}{Harmonic mean}{Weighted arithmetic mean}{}{}}
\LANGpl{
\Question{
Załóżmy, że wartość unikatowego  obrazu  zwiększy się \( 1{,}5 \) razy w pierwszym roku, \( 1{,}2 \) razy w drugim roku \( 1{,}9 \)  razy w trzecim roku. Jesteśmy zainteresowani złożoną roczną stopą wzrostu w tych trzech latach. Jakiego rodzaju średniej potrzebujemy?}
{Średnia geometryczna}{Średnia arytmetyczna}{Średnia harmoniczna}{Ważona średnia arytmetyczna}{}{}}
\LANGcs{
\Question{
Předpokládejme, že vlastníte umělecké dílo, jehož cena vzroste za první rok \( 1{,}5 \) krát, druhý rok \( 1{,}2 \) krát a třetí rok \( 1{,}9 \) krát. Zajímá nás průměrné roční tempo růstu ceny daného uměleckého díla za dané období. Jaký typ průměru se pro výpočet použije?}
{Geometrický průměr}{Aritmetický průměr}{Harmonický průměr}{Vážený aritmetický průměr}{}{}}
\LANGsk{
\Question{
Predpokladajme, že vlastníte umelecké dielo, ktorého cena vzrastie za prvý rok \( 1{,}5 \) krát, druhý rok \( 1{,}2 \) krát a tretí rok \( 1{,}9 \) krát. Zaujíma nás priemerné ročné tempo rastu ceny daného umeleckého diela za dané obdobie. Aký typ priemeru sa pre výpočet použije?}
{Geometrický priemer}{Aritmetický priemer}{Harmonický priemer}{Vážený aritmetický priemer}{}{}}
\LANGes{
\Question{
Suponemos que eres dueño de una obra de arte cuyo precio aumentó \( 1{,}5 \) veces el primer añodruhý rok, el segundo año \( 1{,}2 \) veces y tercer año \( 1{,}9 \) veces. Estamos interesados en su tasa de crecimiento media. ¿Qué tipo de media vamos a calcular?}
{Media geométrica}{Media aritmética}{Media armónica}{Media aritmética ponderada}{}{}}
\End

\Data\ID{1003029507}
\Updated{2021-02-19 01:02:07}
\Level{A}
\SubArea{Statistics}
\LANGen{
\Question{
You recorded the outdoor temperature in the place of your residence at the noontime during the past ten consecutive days. Now, you would like to find out the average noontime outdoor temperature in this place in the ten-day interval. What type of average do you have to use?}
{Arithmetic mean}{Geometric mean}{Harmonic mean}{Weighted geometric mean}{}{}}
\LANGpl{
\Question{
Zarejestrowałeś temperaturę zewnętrzną w miejscu zamieszkania w południe w ciągu ostatnich dziesięciu kolejnych dni. Teraz chcesz poznać średnią temperaturę zewnętrzną w południe w tym miejscu w przedziale dziesięciodniowym. Jakiego rodzaju średniej musisz użyć?}
{Średnia arytmetyczna}{Średnia geometryczna}{Średnia harmoniczna}{Ważona średnia geometryczna}{}{}}
\LANGcs{
\Question{
Deset dní po sobě jste v poledne zapisovali teplotu v místě vašeho bydliště. Nyní chcete zjistit průměrnou polední teplotu v místě vašeho bydliště ve sledovaném období. Jaký typ průměru se pro výpočet použije?}
{Aritmetický průměr}{Geometrický průměr}{Harmonický průměr}{Vážený geometrický průměr}{}{}}
\LANGsk{
\Question{
Desať dní po sebe ste na poludnie zapisovali teplotu v mieste vášho bydliska. Teraz chcete zistiť priemernú poludňajšiu teplotu v mieste vášho bydliska v danom desaťdennom intervale. Aký typ priemeru sa pre výpočet použije?}
{Aritmetický priemer}{Geometrický priemer}{Harmonický priemer}{Vážený geometrický priemer}{}{}}
\LANGes{
\Question{
Durante diez días consecutivos estabas anotando la temperatura en mediodía. Ahora quieres averiguar la temperatura media. ¿Qué tipo de media vas a usar?}
{Media aritmética}{Media geométrica}{Media armónica}{Media aritmética ponderada}{}{}}
\End

\Data\ID{1003029508}
\Updated{2021-02-19 12:02:39}
\Level{A}
\SubArea{Statistics}
\LANGen{
\Question{
There were the following speeds of six cars recorded on one of the roads in New York: $66\,\mathrm{mph}$, $57\,\mathrm{mph}$, $71\,\mathrm{mph}$, $54\,\mathrm{mph}$, $69\,\mathrm{mph}$, $58\,\mathrm{mph}$.
We would like to know the average of these recorded speeds. What type of average do we need to use?}
{Arithmetic mean}{Geometric mean}{Harmonic mean}{Weighted harmonic mean}{}{}}
\LANGpl{
\Question{
Na jednej z dróg w Nowym Jorku zarejestrowano następujące prędkości sześciu samochodów: $66\,\mathrm{km/h}$, $57\,\mathrm{km/h}$, $71\,\mathrm{km/h}$, $54\,\mathrm{km/h}$, $69\,\mathrm{km/h}$, $58\,\mathrm{km/h}$.
Chcielibyśmy poznać średnią z tych zarejestrowanych prędkości. Jakiego rodzaju średniej potrzebujemy?}
{Średnia arytmetyczna}{Średnia geometryczna}{Średnia harmoniczna}{Ważona średnia harmoniczna}{}{}}
\LANGcs{
\Question{
Na cestě v New Yorku byly zaznamenané rychlosti šesti aut: $66\,\mathrm{km/h}$, $57\,\mathrm{km/h}$, $71\,\mathrm{km/h}$, $54\,\mathrm{km/h}$, $69\,\mathrm{km/h}$, $58\,\mathrm{km/h}$.
Zajímá nás průměrná zaznamenaná rychlost těchto šesti aut. Jaký typ průměru se pro výpočet použije?}
{Aritmetický průměr}{Geometrický průměr}{Harmonický průměr}{Vážený harmonický průměr}{}{}}
\LANGsk{
\Question{
Na ceste v New Yorku boli zaznamenané rýchlosti šiestich áut: $66\,\mathrm{km/h}$, $57\,\mathrm{km/h}$, $71\,\mathrm{km/h}$, $54\,\mathrm{km/h}$, $69\,\mathrm{km/h}$, $58\,\mathrm{km/h}$.
Zaujíma nás priemer zaznamenaných rýchlostí týchto aut. Aký typ priemeru sa pre výpočet použije?}
{Aritmetický priemer}{Geometrický priemer}{Harmonický priemer}{Vážený harmonický priemer}{}{}}
\LANGes{
\Question{
Hemos registrado las velocidades de seis coches en una calle en Nueva York: $66\,\mathrm{km/h}$, $57\,\mathrm{km/h}$, $71\,\mathrm{km/h}$, $54\,\mathrm{km/h}$, $69\,\mathrm{km/h}$, $58\,\mathrm{km/h}$.
Estamos interesados en la velocidad media de estos coches. ¿Qué tipo de media vamos a calcular?}
{Media aritmética}{Media geométrica}{Media armónica}{Media aritmética ponderada}{}{}}
\End

\Data\ID{1003123501}
\Updated{2021-03-11 09:03:02}
\Level{A}
\SubArea{Statistics}
\LANGen{
\Question{
At a weather station in Las Vegas wind speeds were measured every day at \( 7 \)pm during one month. The results are specified in the following table.
\[
\begin{array}{|l|c|c|c|c|c|c|c|c|} \hline \text{Day} & 1. & 2. & 3. & 4. & 5. & 6. & 7. & 8. \\\hline \text{Wind} (\mathrm{mps}) & 3 & 2 & 1 & 1 & 2 & 2 & 1 & 3  \\\hline \\\hline  \text{Day} & 9. & 10. & 11. & 12. & 13. & 14. & 15. &  \\\hline  \text{Wind} (\mathrm{mps}) & 2 & 1 & 2 & 2 & 4 & 2 & 4 & \\\hline \\\hline \text{Day} & 16. & 17. & 18. & 19. & 20. & 21. & 22. & 23. \\\hline \text{Wind} (\mathrm{mps}) & 4 & 2 & 2 & 3 & 2 & 12 & 13 & 6  \\\hline \\\hline \text{Day} &  24. & 25. & 26. & 27. & 28. & 29. & 30. &  \\\hline  \text{Wind} (\mathrm{mps}) & 5 & 7 & 2 & 3 & 8 & 9 & 12 & \\\hline\end{array}
\]
Determine the median of the recorded wind speeds.}
{\( 2{.}5\,\mathrm{m/s} \)}{\( 2\,\mathrm{m/s} \)}{\( 4\,\mathrm{m/s} \)}{\( 6\,\mathrm{m/s} \)}{}{}}
\LANGpl{
\Question{
Na stacji meteorologicznej w  Las Vegas mierzono prędkość wiatru codziennie o godzinie \( 19 \) przez miesiąc. Wyniki pomiarów przedstawiono w tabeli.
\[
\begin{array}{|l|c|c|c|c|c|c|c|c|} \hline \text{Dzień} & 1. & 2. & 3. & 4. & 5. & 6. & 7. & 8. \\\hline \text{Wiatr} (\mathrm{m/s}) & 3 & 2 & 1 & 1 & 2 & 2 & 1 & 3  \\\hline \\\hline  \text{Dzień} & 9. & 10. & 11. & 12. & 13. & 14. & 15. &  \\\hline  \text{Wiatr} (\mathrm{m/s}) & 2 & 1 & 2 & 2 & 4 & 2 & 4 & \\\hline \\\hline \text{Dzień} & 16. & 17. & 18. & 19. & 20. & 21. & 22. & 23. \\\hline \text{Wiatr} (\mathrm{m/s}) & 4 & 2 & 2 & 3 & 2 & 12 & 13 & 6  \\\hline \\\hline \text{Dzień} &  24. & 25. & 26. & 27. & 28. & 29. & 30. &  \\\hline  \text{Wiatr} (\mathrm{m/s}) & 5 & 7 & 2 & 3 & 8 & 9 & 12 & \\\hline\end{array}
\]
Oblicz medianę zarejestrowanych pomiarów.}
{\( 2{,}5\,\mathrm{m/s} \)}{\( 2\,\mathrm{m/s} \)}{\( 4\,\mathrm{m/s} \)}{\( 6\,\mathrm{m/s} \)}{}{}}
\LANGcs{
\Question{
Meteorologická stanice v Las Vegas po dobu jednoho měsíce zaznamenávala rychlost větru měřenou v 19:00 hodin. Výsledky jsou uvedeny v následující tabulce:
\[
\begin{array}{|l|c|c|c|c|c|c|c|c|} \hline \text{Den} & 1. & 2. & 3. & 4. & 5. & 6. & 7. & 8. \\\hline \text{Vítr} (\mathrm{m/s}) & 3 & 2 & 1 & 1 & 2 & 2 & 1 & 3  \\\hline \\\hline  \text{Den} & 9. & 10. & 11. & 12. & 13. & 14. & 15. &  \\\hline  \text{Vítr} (\mathrm{m/s}) & 2 & 1 & 2 & 2 & 4 & 2 & 4 & \\\hline \\\hline \text{Den} & 16. & 17. & 18. & 19. & 20. & 21. & 22. & 23. \\\hline \text{Vítr} (\mathrm{m/s}) & 4 & 2 & 2 & 3 & 2 & 12 & 13 & 6  \\\hline \\\hline \text{Den} &  24. & 25. & 26. & 27. & 28. & 29. & 30. & \\\hline  \text{Vítr} (\mathrm{m/s}) & 5 & 7 & 2 & 3 & 8 & 9 & 12 & \\\hline\end{array}
\]
Určete medián zaznamenaných rychlostí větru.}
{\( 2{,}5\,\mathrm{m/s} \)}{\( 2\,\mathrm{m/s} \)}{\( 4\,\mathrm{m/s} \)}{\( 6\,\mathrm{m/s} \)}{}{}}
\LANGsk{
\Question{
Meteorologická stanica v Las Vegas počas jedného mesiaca zaznamenala rýchlosti vetra merané o 19:00 hodine. Výsledky sú uvedené v nasledujúcej tabuľke:
\[
\begin{array}{|l|c|c|c|c|c|c|c|c|} \hline \text{Deň} & 1. & 2. & 3. & 4. & 5. & 6. & 7. & 8. \\\hline \text{Vietor} (\mathrm{m/s}) & 3 & 2 & 1 & 1 & 2 & 2 & 1 & 3  \\\hline \\\hline  \text{Deň} & 9. & 10. & 11. & 12. & 13. & 14. & 15. &  \\\hline \text{Vietor} (\mathrm{m/s}) & 2 & 1 & 2 & 2 & 4 & 2 & 4 & \\\hline \\\hline \text{Deň} & 16. & 17. & 18. & 19. & 20. & 21. & 22. & 23. \\\hline \text{Vietor} (\mathrm{m/s}) & 4 & 2 & 2 & 3 & 2 & 12 & 13 & 6  \\\hline \\\hline \text{Deň} &  24. & 25. & 26. & 27. & 28. & 29. & 30. & \\\hline  \text{Vietor} (\mathrm{m/s}) & 5 & 7 & 2 & 3 & 8 & 9 & 12 & \\\hline\end{array}
\]
Určte medián zaznamenaných rýchlostí vetra.}
{\( 2{,}5\,\mathrm{m/s} \)}{\( 2\,\mathrm{m/s} \)}{\( 4\,\mathrm{m/s} \)}{\( 6\,\mathrm{m/s} \)}{}{}}
\LANGes{
\Question{
La estación meteorológica en Las Vegas anotó la velocidad de viento cada día en 19:00 horas durante un mes. Los resultados están en la tabla:
\[
\begin{array}{|l|c|c|c|c|c|c|c|c|} \hline \text{Día} & 1. & 2. & 3. & 4. & 5. & 6. & 7. & 8. \\\hline \text{Viento} (\mathrm{m/s}) & 3 & 2 & 1 & 1 & 2 & 2 & 1 & 3  \\\hline \\\hline  \text{Día} & 9. & 10. & 11. & 12. & 13. & 14. & 15. &  \\\hline  \text{Viento} (\mathrm{m/s}) & 2 & 1 & 2 & 2 & 4 & 2 & 4 & \\\hline \\\hline \text{Día} & 16. & 17. & 18. & 19. & 20. & 21. & 22. & 23. \\\hline \text{Viento} (\mathrm{m/s}) & 4 & 2 & 2 & 3 & 2 & 12 & 13 & 6  \\\hline \\\hline \text{Día} &  24. & 25. & 26. & 27. & 28. & 29. & 30. & \\\hline  \text{Viento} (\mathrm{m/s}) & 5 & 7 & 2 & 3 & 8 & 9 & 12 & \\\hline\end{array}
\]
Calcula la mediana de la velocidad del viento.}
{\( 2{,}5\,\mathrm{m/s} \)}{\( 2\,\mathrm{m/s} \)}{\( 4\,\mathrm{m/s} \)}{\( 6\,\mathrm{m/s} \)}{}{}}
\End

\Data\ID{1003123502}
\Updated{2021-03-11 09:03:50}
\Level{A}
\SubArea{Statistics}
\LANGen{
\Question{
At a weather station in Las Vegas temperature was measured every day at \( 7 \)pm during one month. The results are specified in the following table.
\[
\begin{array}{|l|c|c|c|c|c|c|c|c|} \hline 
\text{Day} & 1. & 2. & 3. & 4. & 5. & 6. & 7. & 8. \\\hline \text{Temperature }(^{\circ}\mathrm{C}) & 24 & 22 & 21 & 26 & 22 & 23 & 21 & 23 \\\hline \\\hline
\text{Day} & 9. & 10. & 11. & 12. & 13. & 14. & 15. & \\\hline \text{Temperature } (^{\circ}\mathrm{C}) & 21 & 26 & 20 & 23 & 24 & 19 & 21 & \\\hline \\\hline
\text{Day} & 16. & 17. & 18. & 19. & 20. & 21. & 22. & 23.  \\\hline \text{Temperature } (^{\circ}\mathrm{C}) & 21 & 20 & 26 & 23 & 24 & 22 & 23 & 26  \\\hline \\\hline
\text{Day} & 24. & 25. & 26. & 27. & 28. & 29. & 30. &  \\\hline \text{Temperature } (^{\circ}\mathrm{C}) & 25 & 23 & 22 & 25 & 27 & 26 & 22 & \\\hline \end{array}
\]
 Determine the mode of the recorded temperatures.}
{\( 23 \,^{\circ}\mathrm{C} \)}{\( 22\,^{\circ}\mathrm{C} \)}{\( 21\,^{\circ}\mathrm{C} \)}{\( 26\,^{\circ}\mathrm{C} \)}{\( 21\,^{\circ}\mathrm{C} \) also \( 22\,^{\circ}\mathrm{C}\) also \(26\,^{\circ}\mathrm{C}\)}{}}
\LANGpl{
\Question{
Na stacji meteorologicznej w Las Vegas temperaturę mierzono codziennie o godzinie \( 7 \) po południu w ciągu miesiąca. Wyniki są przedstawione  w poniższej tabeli.
\[
\begin{array}{|l|c|c|c|c|c|c|c|c|} \hline 
\text{Dzień} & 1. & 2. & 3. & 4. & 5. & 6. & 7. & 8. \\\hline  \text{Temperatura }(^{\circ}\mathrm{C}) & 24 & 22 & 21 & 26 & 22 & 23 & 21 & 23 \\\hline \\\hline
\text{Dzień} & 9. & 10. & 11. & 12. & 13. & 14. & 15. & \\\hline  \text{Temperatura } (^{\circ}\mathrm{C}) & 21 & 26 & 20 & 23 & 24 & 19 & 21 & \\\hline \\\hline
\text{Dzień} & 16. & 17. & 18. & 19. & 20. & 21. & 22. & 23.  \\\hline  \text{Temperatura } (^{\circ}\mathrm{C}) & 21 & 20 & 26 & 23 & 24 & 22 & 23 & 26  \\\hline \\\hline
\text{Dzień} & 24. & 25. & 26. & 27. & 28. & 29. & 30. &  \\\hline  \text{Temperatura } (^{\circ}\mathrm{C}) & 25 & 23 & 22 & 25 & 27 & 26 & 22 & \\\hline \end{array}
\]
Określ wartość modalną zarejestrowanych temperatur.}
{\( 23 \,^{\circ}\mathrm{C} \)}{\( 22\,^{\circ}\mathrm{C} \)}{\( 21\,^{\circ}\mathrm{C} \)}{\( 26\,^{\circ}\mathrm{C}\)}{\( 21\,^{\circ}\mathrm{C} \) również \( 22\,^{\circ}\mathrm{C}\) również \(26\,^{\circ}\mathrm{C}\)}{}}
\LANGcs{
\Question{
Meteorologická stanice v Las Vegas po dobu jednoho měsíce zaznamenávala teploty měřené v 19:00 hodin. Výsledky jsou uvedeny v následující tabulce:
\[
\begin{array}{|l|c|c|c|c|c|c|c|c|} \hline 
\text{Den} & 1. & 2. & 3. & 4. & 5. & 6. & 7. & 8. \\\hline \text{Teplota } (^{\circ}\mathrm{C}) & 24 & 22 & 21 & 26 & 22 & 23 & 21 & 23 \\\hline \\\hline
\text{Den} & 9. & 10. & 11. & 12. & 13. & 14. & 15. & \\\hline \text{Teplota } (^{\circ}\mathrm{C}) & 21 & 26 & 20 & 23 & 24 & 19 & 21 & \\\hline \\\hline
\text{Den} & 16. & 17. & 18. & 19. & 20. & 21. & 22. & 23.  \\\hline \text{Teplota } (^{\circ}\mathrm{C}) & 21 & 20 & 26 & 23 & 24 & 22 & 23 & 26  \\\hline \\\hline
\text{Den} & 24. & 25. & 26. & 27. & 28. & 29. & 30. &  \\\hline \text{Teplota } (^{\circ}\mathrm{C}) & 25 & 23 & 22 & 25 & 27 & 26 & 22 & \\\hline \end{array}
\]
 Určete modus zaznamenaných teplot.}
{\( 23 \,^{\circ}\mathrm{C} \)}{\( 22\,^{\circ}\mathrm{C} \)}{\( 21\,^{\circ}\mathrm{C} \)}{\( 26\,^{\circ}\mathrm{C}\)}{\( 21\,^{\circ}\mathrm{C} \) i \( 22\,^{\circ}\mathrm{C}\) i \(26\,^{\circ}\mathrm{C}\)}{}}
\LANGsk{
\Question{
Meteorologická stanica v Las Vegas počas jedného mesiaca zaznamenala teploty merané o 19:00 hodine. Výsledky sú uvedené v nasledujúcej tabuľke:
\[
\begin{array}{|l|c|c|c|c|c|c|c|c|} \hline 
\text{Deň} & 1. & 2. & 3. & 4. & 5. & 6. & 7. & 8. \\\hline  \text{Teplota }(^{\circ}\mathrm{C}) & 24 & 22 & 21 & 26 & 22 & 23 & 21 & 23 \\\hline \\\hline
\text{Deň} & 9. & 10. & 11. & 12. & 13. & 14. & 15. & \\\hline  \text{Teplota } (^{\circ}\mathrm{C}) & 21 & 26 & 20 & 23 & 24 & 19 & 21 & \\\hline \\\hline
\text{Deň} & 16. & 17. & 18. & 19. & 20. & 21. & 22. & 23.  \\\hline  \text{Teplota } (^{\circ}\mathrm{C}) & 21 & 20 & 26 & 23 & 24 & 22 & 23 & 26  \\\hline \\\hline
\text{Deň} & 24. & 25. & 26. & 27. & 28. & 29. & 30. &  \\\hline  \text{Teplota } (^{\circ}\mathrm{C}) & 25 & 23 & 22 & 25 & 27 & 26 & 22 & \\\hline \end{array}
\]
 Určte modus zaznamenaných teplôt.}
{\( 23 \,^{\circ}\mathrm{C} \)}{\( 22\,^{\circ}\mathrm{C} \)}{\( 21\,^{\circ}\mathrm{C} \)}{\( 26\,^{\circ}\mathrm{C}\)}{\( 21\,^{\circ}\mathrm{C} \) aj \( 22\,^{\circ}\mathrm{C}\) aj \(26\,^{\circ}\mathrm{C}\)}{}}
\LANGes{
\Question{
La estación meteorológica en Las Vegas anotó la temperatura cada día en 19:00 horas durante un mes. Los resultados están en la tabla:
\[
\begin{array}{|l|c|c|c|c|c|c|c|c|} \hline 
\text{Día} & 1. & 2. & 3. & 4. & 5. & 6. & 7. & 8. \\\hline \text{Temperatura } (^{\circ}\mathrm{C}) & 24 & 22 & 21 & 26 & 22 & 23 & 21 & 23 \\\hline \\\hline
\text{Día} & 9. & 10. & 11. & 12. & 13. & 14. & 15. & \\\hline \text{Temperatura } (^{\circ}\mathrm{C}) & 21 & 26 & 20 & 23 & 24 & 19 & 21 & \\\hline \\\hline
\text{Día} & 16. & 17. & 18. & 19. & 20. & 21. & 22. & 23.  \\\hline \text{Temperatura } (^{\circ}\mathrm{C}) & 21 & 20 & 26 & 23 & 24 & 22 & 23 & 26  \\\hline \\\hline
\text{Día} & 24. & 25. & 26. & 27. & 28. & 29. & 30. &  \\\hline \text{Temperatura } (^{\circ}\mathrm{C}) & 25 & 23 & 22 & 25 & 27 & 26 & 22 & \\\hline \end{array}
\]
Averigua la moda de las temperaturas.}
{\( 23 \,^{\circ}\mathrm{C} \)}{\( 22\,^{\circ}\mathrm{C} \)}{\( 21\,^{\circ}\mathrm{C} \)}{\( 26\,^{\circ}\mathrm{C}\)}{\( 21\,^{\circ}\mathrm{C} \) i \( 22\,^{\circ}\mathrm{C}\) i \(26\,^{\circ}\mathrm{C}\)}{}}
\End

\Data\ID{1003123503}
\Updated{2021-05-04 01:05:53}
\Level{A}
\SubArea{Statistics}
\LANGen{
\Question{
It is known that in the given time period one up to five pieces of the same product were sold in each of the 100 monitored shops. One piece was sold in \( 26 \) shops, \( 2 \) pieces in \( 64 \) shops, \( 3 \) pieces in \( 7 \) shops, \( 4 \) pieces in \( 2 \) shops and \( 5 \) pieces in one shop. How many pieces of the product  were  sold  most frequently  in the given shops? Choose the correct characteristic and its value.}
{Mode: \( 2 \) pieces}{Arithmetic mean: \( 3 \) pieces}{Median: \( 2 \) pieces}{Median: \( 3 \) pieces}{(Weighted) arithmetic mean: \( 1.88 \) pieces}{}}
\LANGpl{
\Question{
Wiadomo, że  w określonym czasie sprzedano od jednej do pięciu sztuk tego samego produktu  w każdym ze  100 monitorowanych sklepów. Jedną sztukę sprzedano w  \( 26 \) sklepach, \( 2 \) sztuki w  \( 64 \) sklepach, \( 3 \) sztuki w  \( 7 \) sklepach, \( 4 \) sztuki w  \( 2 \) sklepach oraz \( 5 \) sztuk w jednym sklepie. Ile sztuk tego produktu było sprzedawanych najczęściej w podanych sklepach. Wybierz poprawną charakterystykę i jej wartość.}
{Dominanta: \( 2 \) sztuki}{Średnia arytmetyczna: \( 3 \) sztuki}{Mediana: \( 2 \) sztuki}{Mediana: \( 3 \) sztuki}{(Ważona ) średnia arytmetyczna: \( 1{,}88 \) sztuki}{}}
\LANGcs{
\Question{
Víme, že v daném období v každém ze \( 100 \) sledovaných obchodů prodali jeden až pět kusů určitého zboží. Jeden kus se prodal v \( 26 \) obchodech, \( 2 \) kusy v \( 64 \) obchodech, \( 3 \) kusy v \( 7 \) obchodech, \( 4 \) kusy ve \( 2 \) obchodech a \( 5 \) kusů v jednom obchodě. V jakém množství se v daných obchodech toto zboží prodávalo nejčastěji? Vyberte správnou charakteristiku a její hodnotu.}
{Modus: \( 2 \) kusy}{Aritmetický průměr: \( 3 \) kusy}{Medián: \( 2 \) kusy}{Medián: \( 3 \) kusy}{(Vážený) aritmetický průměr: \( 1{,}88 \) kusů}{}}
\LANGsk{
\Question{
Je známe, že v danom období sa v každom zo \( 100 \) sledovaných obchodov predalo jeden až päť kusov istého tovaru. Jeden kus sa predal v \( 26 \) obchodoch, \( 2 \) kusy v \( 64 \) obchodoch, \( 3 \) kusy v \( 7 \) obchodoch, \( 4 \) kusy v \( 2 \) obchodoch a \( 5 \) kusov v jednom obchode. V akom množstve sa v daných obchodoch tento tovar predával najčastejšie? Vyberte správnu charakteristiku a jej hodnotu.}
{Modus: \( 2 \) kusy}{Aritmetický priemer: \( 3 \) kusy}{Medián: \( 2 \) kusy}{Medián: \( 3 \) kusy}{(Vážený) aritmetický priemer: \( 1{,}88 \) kusov}{}}
\LANGes{
\Question{
Sabemos que durante un periodo monitorado se vendieron de uno hasta cinco piezas de un tipo de producto en uno de las \( 100 \) tiendas monitoradas. En \( 26 \) tiendas vendieron una pieza, dos piezas en \( 64 \) tiendas, tres piezas en \( 7 \) tiendas, cuatro piezas en \( 2 \) riendas y cinco piezas en una tienda. ¿Qué cantidad de productos era más frecuente? Elije la característica y su valor.}
{Moda: \( 2 \) piezas}{Media aritmética: \( 3 \) priezas}{Mediana: \( 2 \) piezas}{Mediana: \( 3 \) piezas}{Media aritmética valorada: \( 1{,}88 \) piezas}{}}
\End

\Data\ID{1103025101}
\Updated{2020-07-15 03:07:12}
\Level{A}
\SubArea{Statistics}
\IMG{1103025101.pdf}
\LANGen{
\Question{
The results of the math test are shown in the graph. Based on the graph, which of the following statements is false?}
{The median of the scores is the same as their modus.}{Half of the students had higher score than the average score.}{The average score rounded to the nearest hundredth is \( 2{.}68 \).}{}{}{}}
\LANGpl{
\Question{
Wyniki testu matematycznego pokazano na wykresie. Na podstawie wykresu, wskaż  stwierdzenie  fałszywe.}
{Mediana wyników jest taka sama jak ich modus.}{Połowa uczniów miała wyższy wynik niż średni wynik.}{Średni wynik zaokrąglony do dwóch miejsc po przecinku wynosi \( 2{,}68 \).}{}{}{}}
\LANGcs{
\Question{
Na obrázku jsou znázorněné výsledky písemné práce z matematiky. Určete, které z daných tvrzení o příslušném statistickém souboru je nepravdivé.
(Slovníček: Number of students - Počet studentů, Score - Známka)}
{Medián známek je stejný jako jejich modus.}{Polovina žáků dostala z této písemné práce horší než průměrnou známku.}{Průměrná známka z této písemné práce, vypočítaná s přesností na dvě desetinná místa, je \( 2{,}68 \).}{}{}{}}
\LANGsk{
\Question{
Na obrázku sú znázornené výsledky písomnej práce z matematiky. Určite, ktoré z daných tvrdení o príslušnom štatistickom súbore je nepravdivé.
(Slovníček: Number of students - Počet študentov, Score - Známka)}
{Medián známok je taký istý ako ich modus.}{Polovica žiakov dostala z tejto písomnej práce horšiu ako priemernú známku.}{Priemerná známka z tejto písomnej práce, vypočítaná s presnosťou na dve desatinné miesta, je \( 2{,}68 \).}{}{}{}}
\LANGes{
\Question{
The results of the math test are shown in the graph. Based on the graph, which of the following statements is false?}
{The median of the scores is the same as their modus.}{Half of the students had higher score than the average score.}{The average score rounded to the nearest hundredth is \( 2{.}68 \).}{}{}{}}
\End

\Data\ID{9000139501}
\Updated{2020-07-15 03:07:37}
\Level{A}
\SubArea{Statistics}
\LANGen{
\Question{
Ten apples in a box have average mass
\(200\, \mathrm{g}\). We remove one
apple of the mass \(200\, \mathrm{g}\) 
from the box. What is the change in the average mass of the apples from the box?}
{The average mass of the apples does not change.}{The average mass of the apples decreases by
\(20\, \mathrm{g}\).}{The average mass of the apples increases by
\(20\, \mathrm{g}\).}{There is not enough information to solve this problem.}{}{}}
\LANGpl{
\Question{
Średnia waga dziesięciu jabłek znajdujących się w pudełku wynosi
\(200\, \mathrm{g}\). Wyciągamy jedno jabłko z pudełka o  wadze \(200\, \mathrm{g}\).
Jak zmieni się średnia waga jabłek w pudełku?}
{Nie zmieni się.}{Zmniejszy się o
\(20\, \mathrm{g}\).}{Zwiększy się o 
\(20\, \mathrm{g}\).}{Za mało informacji, aby udzielić odpowiedzi.}{}{}}
\LANGcs{
\Question{
Průměrná hmotnost deseti jablek činí
\(200\, \mathrm{g}\). Jak
se změní průměrná hmotnost jablek, jestliže jedno jablko o hmotnosti
\(200\, \mathrm{g}\) 
sníme?}
{Nezmění se.}{Klesne o \(20\, \mathrm{g}\).}{Vzroste o \(20\, \mathrm{g}\).}{Nelze určit.}{}{}}
\LANGsk{
\Question{
Priemerná hmotnosť desiatich jabĺk v krabici je
\(200\, \mathrm{g}\). Ako sa zmení priemerná hmotnosť jabĺk, ak jedno jablko o hmotnosti \(200\, \mathrm{g}\) 
odstránime z krabice?}
{Priemerná hmotnosť jabĺk sa nezmení.}{Priemerná hmotnosť jabĺk klesne o
\(20\, \mathrm{g}\).}{Priemerná hmotnosť jabĺk vzrastie o 
\(20\, \mathrm{g}\).}{Nedá sa určiť.}{}{}}
\LANGes{
\Question{
Ten apples in a box have average mass
\(200\, \mathrm{g}\). We remove one
apple of the mass \(200\, \mathrm{g}\) 
from the box. What is the change in the average mass of the apples from the box?}
{The average mass of the apples does not change.}{The average mass of the apples decreases by
\(20\, \mathrm{g}\).}{The average mass of the apples increases by
\(20\, \mathrm{g}\).}{There is not enough information to solve this problem.}{}{}}
\End

\Data\ID{9000139502}
\Updated{2020-07-15 03:07:37}
\Level{A}
\SubArea{Statistics}
\LANGen{
\Question{
The average mass of \(30\) 
eggs on a plate is \(60\, \mathrm{g}\).
From this amount we remove five eggs. The total mass of these five eggs is
\(280\, \mathrm{g}\).
Find the change in the average mass of the remaining eggs on the plate.}
{The average mass of eggs increases by \(0.8\, \mathrm{g}\).}{The average mass of eggs decreases by
\(4\, \mathrm{g}\).}{The average mass of eggs increases by
\(4\, \mathrm{g}\).}{The average mass of eggs increases by
\(12\, \mathrm{g}\).}{}{}}
\LANGpl{
\Question{
Średnia waga  \(30\) 
jajek na talerzu wynosi \(60\, \mathrm{g}\).
Zabieramy pięć jajek o całkowitej wadze
\(280\, \mathrm{g}\).
Jak zmieniła się średnia waga pozostałych jajek na talerzu?}
{Zwiększyła się o  \(0{,}8\, \mathrm{g}\).}{Zmniejszyła się o
\(4\, \mathrm{g}\).}{Zwiększyła się o 
\(4\, \mathrm{g}\).}{Zwiększyła  się o
\(12\, \mathrm{g}\).}{}{}}
\LANGcs{
\Question{
Průměrná hmotnost třiceti vajec činí
\(60\, \mathrm{g}\). Jak se
změní tato průměrná hmotnost, jestliže z pěti vajec o celkové hmotnosti
\(280\, \mathrm{g}\) 
uděláme omeletu?}
{Vzroste o \(0{,}8\, \mathrm{g}\).}{Klesne o \(4\, \mathrm{g}\).}{Vzroste o \(4\, \mathrm{g}\).}{Vzroste o \(12\, \mathrm{g}\).}{}{}}
\LANGsk{
\Question{
Priemerná hmotnosť \(30\) 
vajec na tanieri je \(60\, \mathrm{g}\).
Ako sa zmení táto priemerná hmotnosť, ak vezmeme z taniera päť vajec o hmotnosti
\(280\, \mathrm{g}\)?}
{Priemerná hmotnosť vajec vzrastie o \(0{,}8\, \mathrm{g}\).}{Priemerná hmotnosť vajec klesne o
\(4\, \mathrm{g}\).}{Priemerná hmotnosť vajec vzrastie o
\(4\, \mathrm{g}\).}{Priemerná hmotnosť vajec vzrastie o
\(12\, \mathrm{g}\).}{}{}}
\LANGes{
\Question{
The average mass of \(30\) 
eggs on a plate is \(60\, \mathrm{g}\).
From this amount we remove five eggs. The total mass of these five eggs is
\(280\, \mathrm{g}\).
Find the change in the average mass of the remaining eggs on the plate.}
{The average mass of eggs increases by \(0.8\, \mathrm{g}\).}{The average mass of eggs decreases by
\(4\, \mathrm{g}\).}{The average mass of eggs increases by
\(4\, \mathrm{g}\).}{The average mass of eggs increases by
\(12\, \mathrm{g}\).}{}{}}
\End

\Data\ID{9000139503}
\Updated{2020-07-15 03:07:37}
\Level{A}
\SubArea{Statistics}
\LANGen{
\Question{
The average mass of the pear in a basket is
\(150\, \mathrm{g}\). Find
the change in the average mass of the pears in the basket if one pear has been
removed from the basket.}
{There is not enough information to solve this problem.}{The average mass of the pears increases by
\(7.5\, \mathrm{g}\).}{The average mass of the pears decreases by
\(7.5\, \mathrm{g}\).}{The average mass of the pears does not change.}{}{}}
\LANGpl{
\Question{
Średnia waga  gruszki w koszyku wynosi 
\(150\, \mathrm{g}\).  W jaki sposób zmieni się średnia waga gruszek w koszyku, jeżeli jedna gruszka zostanie usunięta z koszyka?}
{Za mało informacji, aby udzielić odpowiedzi.}{Zwiększy się o
\(7{,}5\, \mathrm{g}\).}{Zmniejszy się o
\(7{,}5\, \mathrm{g}\).}{Nie zmieni się.}{}{}}
\LANGcs{
\Question{
Průměrná hmotnost dvaceti hrušek činí
\(150\, \mathrm{g}\). Jak
se změní průměrná hmotnost hrušek, jestliže jednu hrušku sníme?}
{Nelze určit.}{Klesne o \(7{,}5\, \mathrm{g}\).}{Vzroste o \(7{,}5\, \mathrm{g}\).}{Nezmění se.}{}{}}
\LANGsk{
\Question{
Priemerná hmotnosť hrušky v košíku je 
\(150\, \mathrm{g}\). Ako sa zmení priemerná hmotnosť hrušiek, ak z košíka vezmeme jednu hrušku?}
{Nedá sa určiť.}{Priemerná hmotnosť vzrastie o 
\(7{,}5\, \mathrm{g}\).}{Priemerná hmotnosť hrušiek klesne o 
\(7{,}5\, \mathrm{g}\).}{Priemerná hmotnosť hrušiek sa nezmení.}{}{}}
\LANGes{
\Question{
The average mass of the pear in a basket is
\(150\, \mathrm{g}\). Find
the change in the average mass of the pears in the basket if one pear has been
removed from the basket.}
{There is not enough information to solve this problem.}{The average mass of the pears increases by
\(7.5\, \mathrm{g}\).}{The average mass of the pears decreases by
\(7.5\, \mathrm{g}\).}{The average mass of the pears does not change.}{}{}}
\End

\Data\ID{9000139504}
\Updated{2020-07-15 03:07:37}
\Level{A}
\SubArea{Statistics}
\LANGen{
\Question{
The average salary of five employees is
\(3\: 000\, \mathrm{Euro}\). This
group of the employees is expanded by one new person. The salary of the new person
is \(2\: 400\, \mathrm{Euro}\).
Find the change in the average salary of this group.}
{The average salary decreases by \(100\, \mathrm{Euro}\).}{The average salary decreases by \(480\, \mathrm{Euro}\).}{The average salary increases by \(400\, \mathrm{Euro}\).}{The average salary increases by \(480\, \mathrm{Euro}\).}{}{}}
\LANGpl{
\Question{
Średnia pensja grupy pięciu pracowników wynosi
\(3\: 000\, \mathrm{Euro}\). Grupę tę powiększono o jedną osobę, jej pensja wynosi 
 \(2\: 400\, \mathrm{Euro}\).
Jak zmieni się średnia pensja pracowników?}
{Zmniejszy się o  \(100\, \mathrm{Euro}\).}{Zmniejszy się o \(480\, \mathrm{Euro}\).}{Zwiększy się o  \(400\, \mathrm{Euro}\).}{Zwiększy się o  \(480\, \mathrm{Euro}\).}{}{}}
\LANGcs{
\Question{
Průměrná odměna pěti zaměstnanců v oddělení činí
\(3\: 000\, \mathrm{K\check{c}}\). Jak
se změní výše průměrné odměny u zaměstnanců v oddělení, jestliže
do jejich pracovního kolektivu přijde nový pracovník, který dostane odměnu
\(2\: 400\, \mathrm{K\check{c}}\)?}
{Klesne o \(100\, \mathrm{K\check{c}}\).}{Klesne o \(480\, \mathrm{K\check{c}}\).}{Vzroste o \(400\, \mathrm{K\check{c}}\).}{Vzroste o \(480\, \mathrm{K\check{c}}\).}{}{}}
\LANGsk{
\Question{
Priemerný plat piatich zamestnancov na oddelení je
\(3\: 000\, \mathrm{eur}\). Ako sa zmení výška priemerného platu u zamestnancov  na oddelení, ak do pracovného kolektívu príjmu nového pracovníka, ktorý dostane plat
 \(2\: 400\, \mathrm{eur}\)?}
{Priemerný plat klesne o \(100\, \mathrm{eur}\).}{Priemerný plat klesne o \(480\, \mathrm{eur}\).}{Priemerný plat narastie o \(400\, \mathrm{eur}\).}{Priemerný plat narastie o \(480\, \mathrm{eur}\).}{}{}}
\LANGes{
\Question{
The average salary of five employees is
\(3\: 000\, \mathrm{Euro}\). This
group of the employees is expanded by one new person. The salary of the new person
is \(2\: 400\, \mathrm{Euro}\).
Find the change in the average salary of this group.}
{The average salary decreases by \(100\, \mathrm{Euro}\).}{The average salary decreases by \(480\, \mathrm{Euro}\).}{The average salary increases by \(400\, \mathrm{Euro}\).}{The average salary increases by \(480\, \mathrm{Euro}\).}{}{}}
\End

\Data\ID{9000139505}
\Updated{2020-07-15 03:07:37}
\Level{A}
\SubArea{Statistics}
\LANGen{
\Question{
The average mass of twelve oranges is
\(120\, \mathrm{g}\).
To this amount we add another six oranges with the average mass
\(150\, \mathrm{g}\). Find
the change in the average mass of oranges.}
{The average mass increases by \(10\, \mathrm{g}\).}{The average mass increases by \(8.3\, \mathrm{g}\).}{The average mass increases by \(25\, \mathrm{g}\).}{The average mass decreases by \(8.3\, \mathrm{g}\).}{}{}}
\LANGpl{
\Question{
Średnia waga dwunastu pomarańczy wynosi 
\(120\, \mathrm{g}\).
Dodajemy jeszcze sześć pomarańczy o średniej wadze 
\(150\, \mathrm{g}\). Jak zmieni się średnia waga pomarańczy?}
{Zwiększy się o  \(10\, \mathrm{g}\).}{Zwiększy się o  \(8{,}3\, \mathrm{g}\).}{Zwiększy się o  \(25\, \mathrm{g}\).}{Zmniejszy się o \(8{,}3\, \mathrm{g}\).}{}{}}
\LANGcs{
\Question{
Průměrná hmotnost dvanácti pomerančů činí
\(120\, \mathrm{g}\).
Jak se změní průměrná hmotnost pomerančů, jestliže k nim
přidáme dalších šest pomerančů s průměrnou hmotností
\(150\, \mathrm{g}\)?}
{Vzroste o \(10\, \mathrm{g}\).}{Vzroste o \(8{,}3\, \mathrm{g}\).}{Vzroste o \(25\, \mathrm{g}\).}{Klesne o \(8{,}3\, \mathrm{g}\).}{}{}}
\LANGsk{
\Question{
Priemerná hmotnosť dvanástich pomarančov je
\(120\, \mathrm{g}\).
Ako sa zmení táto priemerná hmotnosť, ak k nim pridáme ďalších šesť pomarančov s priemernou hmotnosťou
\(150\, \mathrm{g}\)?}
{Vzrastie o \(10\, \mathrm{g}\).}{Vzrastie o  \(8{,}3\, \mathrm{g}\).}{Vzrastie o \(25\, \mathrm{g}\).}{Klesne o \(8{,}3\, \mathrm{g}\).}{}{}}
\LANGes{
\Question{
The average mass of twelve oranges is
\(120\, \mathrm{g}\).
To this amount we add another six oranges with the average mass
\(150\, \mathrm{g}\). Find
the change in the average mass of oranges.}
{The average mass increases by \(10\, \mathrm{g}\).}{The average mass increases by \(8.3\, \mathrm{g}\).}{The average mass increases by \(25\, \mathrm{g}\).}{The average mass decreases by \(8.3\, \mathrm{g}\).}{}{}}
\End

\Data\ID{9000139506}
\Updated{2020-07-15 03:07:37}
\Level{A}
\SubArea{Statistics}
\LANGen{
\Question{
There are eight mandarins of average mass
\(90\, \mathrm{g}\) in the box. We
got two another mandarins and add them to the box. The new average mass of the mandarins
in the box is \(92\, \mathrm{g}\).
Find the average mass of the two added mandarins.}
{\(100\, \mathrm{g}\)}{\(92\, \mathrm{g}\)}{\(96\, \mathrm{g}\)}{\(106\, \mathrm{g}\)}{}{}}
\LANGpl{
\Question{
W koszyku znajduje się osiem mandarynek o średniej wadze 
\(90\, \mathrm{g}\). Do koszyka dodano dwie mandarynki. Po dodaniu dwóch mandarynek średnia waga   mandarynek wynosi
\(92\, \mathrm{g}\).
Ile wynosi średnia waga dwóch dodanych mandarynek?}
{\(100\, \mathrm{g}\)}{\(92\, \mathrm{g}\)}{\(96\, \mathrm{g}\)}{\(106\, \mathrm{g}\)}{}{}}
\LANGcs{
\Question{
Průměrná hmotnost osmi mandarinek činí
\(90\, \mathrm{g}\).
Náhodně vybereme další dvě mandarinky, které přidáme k
původním. Jaká je průměrná hmotnost těchto dvou přidaných
mandarinek, jestliže průměrná hmotnost všech mandarinek vzroste na
\(92\, \mathrm{g}\)?}
{\(100\, \mathrm{g}\)}{\(92\, \mathrm{g}\)}{\(96\, \mathrm{g}\)}{\(106\, \mathrm{g}\)}{}{}}
\LANGsk{
\Question{
Priemerná hmotnosť ôsmych mandaríniek je
\(90\, \mathrm{g}\). Náhodne vyberieme ďalšie dve mandarínky, ktoré pridáme k pôvodným. Aká je priemerná hmotnosť týchto dvoch pridaných mandaríniek, ak priemerná hmotnosť všetkých mandaríniek narastie na \(92\, \mathrm{g}\)?}
{\(100\, \mathrm{g}\)}{\(92\, \mathrm{g}\)}{\(96\, \mathrm{g}\)}{\(106\, \mathrm{g}\)}{}{}}
\LANGes{
\Question{
There are eight mandarins of average mass
\(90\, \mathrm{g}\) in the box. We
got two another mandarins and add them to the box. The new average mass of the mandarins
in the box is \(92\, \mathrm{g}\).
Find the average mass of the two added mandarins.}
{\(100\, \mathrm{g}\)}{\(92\, \mathrm{g}\)}{\(96\, \mathrm{g}\)}{\(106\, \mathrm{g}\)}{}{}}
\End

\Data\ID{9000139507}
\Updated{2021-06-15 09:06:35}
\Level{A}
\SubArea{Statistics}
\LANGen{
\Question{
The average mass of five melons is \(2\: 400\, \mathrm{g}\).
We have to add another melon such that the new average value of all six melons will
be \(2\: 420\, \mathrm{g}\).
Find the mass of the sixth melon.}
{\(2\: 520\, \mathrm{g}\)}{\(2\: 540\, \mathrm{g}\)}{\(2\: 480\, \mathrm{g}\)}{\(2\: 460\, \mathrm{g}\)}{}{}}
\LANGpl{
\Question{
Średnia waga pięciu melonów wynosi \(2\: 400\, \mathrm{g}\).
Musimy dodać kolejnego melona tak, aby  średnia waga sześciu melonów wynosiła 
 \(2\: 420\, \mathrm{g}\).
Wskaż wagę szóstego melona.}
{\(2\: 520\, \mathrm{g}\)}{\(2\: 540\, \mathrm{g}\)}{\(2\: 480\, \mathrm{g}\)}{\(2\: 460\, \mathrm{g}\)}{}{}}
\LANGcs{
\Question{
Průměrná hmotnost pěti melounů činí
\(2\: 400\, \mathrm{g}\).
Určete hmotnost melounu, který musíme k těmto pěti
přidat, aby průměrná hmotnost všech melounů byla
\(2\: 420\, \mathrm{g}\).}
{\(2\: 520\, \mathrm{g}\)}{\(2\: 540\, \mathrm{g}\)}{\(2\: 480\, \mathrm{g}\)}{\(2\: 460\, \mathrm{g}\)}{}{}}
\LANGsk{
\Question{
Priemerná hmotnosť piatich melónov je \(2\: 400\, \mathrm{g}\). Určte hmotnosť melóna, ktorý musíme k týmto piatim pridať, aby priemerná hmotnosť všetkých šiestich melónov bola \(2\: 420\, \mathrm{g}\).}
{\(2\: 520\, \mathrm{g}\)}{\(2\: 540\, \mathrm{g}\)}{\(2\: 480\, \mathrm{g}\)}{\(2\: 460\, \mathrm{g}\)}{}{}}
\LANGes{
\Question{
La masa media de cinco sandías es 
\(2\: 400\, \mathrm{g}\).
Calcula la masa de la sexta sandía con la cuál la masa media será
\(2\: 420\, \mathrm{g}\).}
{\(2\: 520\, \mathrm{g}\)}{\(2\: 540\, \mathrm{g}\)}{\(2\: 480\, \mathrm{g}\)}{\(2\: 460\, \mathrm{g}\)}{}{}}
\End

\Data\ID{9000139509}
\Updated{2021-06-15 09:06:18}
\Level{A}
\SubArea{Statistics}
\LANGen{
\Question{
The average year income in a company was
\(200\: 000\, \mathrm{Euro}\) two years ago. This
income increased by \(10\%\) 
a year ago and by \(80\: 000\, \mathrm{Euro}\) 
this year. Find the average percentage increase in the average income per one year
and round to the nearest percent.}
{\(22\%\)}{\(23\%\)}{\(25\%\)}{\(50\%\)}{}{}}
\LANGpl{
\Question{
Dwa lata temu średni dochód firmy wynosił 
\(200\: 000\, \mathrm{Euro}\). Rok temu dochód zwiększono o
 \(10\%\) 
a w roku obecnym o \(80\: 000\, \mathrm{Euro}\).
Jaki jest średni roczny procentowy wzrost dochodu w podanym okresie?
Zaokrągli odpowiedź do pełnych procentów.}
{\(22\%\)}{\(23\%\)}{\(25\%\)}{\(50\%\)}{}{}}
\LANGcs{
\Question{
Předloni byla výše ročního platu zaměstnance ve firmě
\(200\: 000\, \mathrm{K\check{c}}\), loni vzrostla o
\(10\:\%\) a letos byl roční
plat zaměstnance o \(80\: 000\, \mathrm{K\check{c}}\) 
vyšší než loni. Jaký je průměrný roční procentuální nárůst jeho
platu za sledované období? (zaokrouhlete na procenta)}
{\(22\:\%\)}{\(23\:\%\)}{\(25\:\%\)}{\(50\:\%\)}{}{}}
\LANGsk{
\Question{
Pred dvoma rokmi bola výška ročného platu zamestnanca vo firme 
\(200\: 000\, \mathrm{eur}\), minulý rok vzrástla o \(10\%\) a tento rok o \(80\: 000\, \mathrm{eur}\). Aký je priemerný ročný percentuálny nárast jeho platu za sledované obdobie? (Zaokrúhlite na percentá.)}
{\(22\%\)}{\(23\%\)}{\(25\%\)}{\(50\%\)}{}{}}
\LANGes{
\Question{
El salario de un empleado de un año hace dos años fue
\(200\: 000\, \mathrm{K\check{c}}\), el año pasado creció por
\(10\:\%\) y este año ha sido el salario de un empleado de un año \(80\: 000\, \mathrm{K\check{c}}\) 
más que el año pasado. Averigua el crecimiento anual medio  en el periodo mencionado. (aproxima a porcentajes)}
{\(22\:\%\)}{\(23\:\%\)}{\(25\:\%\)}{\(50\:\%\)}{}{}}
\End

\Data\ID{9000139510}
\Updated{2021-06-15 08:06:24}
\Level{A}
\SubArea{Statistics}
\LANGen{
\Question{
The price of a butter increased by \(8\%\) 
in the year \(2013\) 
and by \(34\%\) in
the year \(2014\).
Find the average percentage growth of the price of the butter per one year in the period
\(2012\)-\(2014\).
Round your answer to the nearest percent.}
{\(20\%\)}{\(21\%\)}{\(14\%\)}{\(26\%\)}{}{}}
\LANGpl{
\Question{
Cena masła wzrosła o \(8\%\) 
 w \(2013\), o 
 \(34\%\) w
 \(2014\) roku.
Ile wynosi roczny wzrost ceny masła w latach 
\(2012\)-\(2014\)?
Zaokrągli odpowiedź do pełnych procentów.}
{\(20\%\)}{\(21\%\)}{\(14\%\)}{\(26\%\)}{}{}}
\LANGcs{
\Question{
V roce \(2013\) byl meziroční
nárůst cen másla \(8\:\%\),
v roce \(2014\) byl meziroční
nárůst cen másla \(34\:\%\).
Jaký byl průměrný meziroční nárůst cen másla v letech
\(2012\) až
\(2014\)?
(zaokrouhlete na procenta)}
{\(20\:\%\)}{\(21\:\%\)}{\(14\:\%\)}{\(26\:\%\)}{}{}}
\LANGsk{
\Question{
Medziročný nárast cien masla bol \(8\%\) 
v roku \(2013\) 
a \(34\%\) v roku \(2014\).
Aký bol priemerný medziročný nárast cien masla v rokoch
\(2012\)-\(2014\)?
Zaokrúhlite na percentá.}
{\(20\%\)}{\(21\%\)}{\(14\%\)}{\(26\%\)}{}{}}
\LANGes{
\Question{
En el año \(2013\) el crecimiento interanual del precio de mantequilla fue \(8\:\%\) y
en el año \(2014\)  fue \(34\:\%\).
¿Cuál fue el crecimiento interanual promedio entre los años  
\(2012\) y
\(2014\)?
(aproxima a porcentajes)}
{\(20\:\%\)}{\(21\:\%\)}{\(14\:\%\)}{\(26\:\%\)}{}{}}
\End

\Data\ID{1003025201}
\Updated{2021-08-23 02:08:19}
\Level{B}
\SubArea{Statistics}
\LANGen{
\Question{
Two hunters, Adam and Boris, competed in target shooting. Adam hit the target points \( \{10;10;9;8;7\}\), and Boris \( \{10;10;9;9;6\} \). Who is the winner? In the case of the same sum of gained points the shooting accuracy is decisive. Which of the following statements is true, if the accuracy is quantified by the variance of the points? (The variance is rounded to two decimal places.)}
{Adam won with the variance of \( 1{.}36\,\mathrm{points}^2 \).}{Adam won with the variance of \( 1{.}17\,\mathrm{points}^2 \).}{Boris won with the variance of \( 2{.}16\,\mathrm{points}^2 \).}{Adam won with the variance of \( 1{.}36\,\mathrm{points} \).}{Adam won with the variance of \( 1{.}17\,\mathrm{points} \).}{Boris won with the variance of \( 2{.}16\,\mathrm{points} \).}}
\LANGpl{
\Question{
Dwóch myśliwych, Adam i Borys strzelało do celu. Adam zdobył  \( \{10;10;9;8;7\}\), a Borys  \( \{10;10;9;9;6\} \). Kto jest zwycięzcą? W przypadku tej samej ilości zdobytych punktów decydująca jest dokładność strzelania. Które z poniższych stwierdzeń jest prawdziwe, jeśli dokładność jest określona wariancją punktów? (Wariancja jest zaokrąglona do dwóch miejsc po przecinku.)}
{Adam wygrał wariancją równą  \( 1{,}36\,\mathrm{pkt}^2 \).}{Adam wygrał wariancją równą \( 1{,}17\,\mathrm{pkt}^2 \).}{Borys wygrał wariancją równą  \( 2{,}16\,\mathrm{pkt}^2 \).}{Adam wygrał wariancją równą  \( 1{,}36\,\mathrm{pkt} \).}{Adam wygrał wariancją równą  \( 1{,}17\,\mathrm{pkt} \).}{Borys wygrał wariancją równą \( 2{,}16\,\mathrm{pkt} \).}}
\LANGcs{
\Question{
Dva myslivci, Adam a Boris soutěžili ve střelbě na terč. Adam trefil v terči body \( \{10;10;9;8;7\}\) a Boris \( \{10;10;9;9;6\} \). Který z nich soutěž vyhrál, jestliže v případě stejného součtu bodů rozhoduje přesnost střelby, tedy bodový rozptyl zásahů?  (Rozptyl zaokrouhlete na dvě desetinná místa.)}
{Vyhrál Adam s rozptylem \( 1{,}36 \) bodů.}{Vyhrál Adam s rozptylem  \( 1{,}17 \) bodů.}{Vyhrál Boris s rozptylem  \( 2{,}16 \) bodů.}{Vyhrál Adam s rozptylem \( 1{,}36 \) bodů.}{Vyhrál Adam s rozptylem  \( 1{,}17 \) bodů.}{Vyhrál Boris s rozptylem  \( 2{,}16 \) bodů.}}
\LANGsk{
\Question{
Dvaja poľovníci, Adam a Boris súťažili v streľbe na terč. Adam trafil v terči body \( \{10;10;9;8;7\}\) a Boris \( \{10;10;9;9;6\} \). Ktorý z nich súťaž vyhral, ak v prípade rovnakého súčtu bodov rozhoduje presnosť streľby, teda bodový rozptyl zásahov?  (Rozptyl zaokrúhlite na dve desatinné miesta.)}
{Vyhral Adam s rozptylom \( 1{,}36\,\mathrm{bodov}^2 \).}{Vyhral Adam s rozptylom  \( 1{,}17\,\mathrm{bodov}^2 \).}{Vyhral Boris s rozptylom  \( 2{,}16\,\mathrm{bodov}^2 \).}{Vyhral Adam s rozptylom \( 1{,}36\,\mathrm{bodov} \).}{Vyhral Adam s rozptylom  \( 1{,}17\,\mathrm{bodov} \).}{Vyhral Boris s rozptylom  \( 2{,}16\,\mathrm{bodov} \).}}
\LANGes{
\Question{
Two hunters, Adam and Boris, competed in target shooting. Adam hit the target points \( \{10;10;9;8;7\}\), and Boris \( \{10;10;9;9;6\} \). Who is the winner? In the case of the same sum of gained points the shooting accuracy is decisive. Which of the following statements is true, if the accuracy is quantified by the variance of the points? (The variance is rounded to two decimal places.)}
{Adam won with the variance of \( 1{.}36\,\mathrm{points}^2 \).}{Adam won with the variance of \( 1{.}17\,\mathrm{points}^2 \).}{Boris won with the variance of \( 2{.}16\,\mathrm{points}^2 \).}{Adam won with the variance of \( 1{.}36\,\mathrm{points} \).}{Adam won with the variance of \( 1{.}17\,\mathrm{points} \).}{Boris won with the variance of \( 2{.}16\,\mathrm{points} \).}}
\End

\Data\ID{1003029401}
\Updated{2020-07-15 03:07:12}
\Level{B}
\SubArea{Statistics}
\LANGen{
\Question{
Wooden boards were expected to be cut to the same length. After their cutting the resulting measured lengths are: \( 2{.}00;\ 2{.}02;\ 2{.}05;\ 2{.}02;\ 2{.}08;\ 2{.}11 \) (in meters). We use the standard deviation of board lengths to describe (quantify) the accuracy of cutting. Find the standard deviation of board lengths rounded to four decimal places.}
{\( 0{.}0382\,\mathrm{m} \)}{\( 0{.}0381\,\mathrm{m} \)}{\( 0{.}0014\,\mathrm{m} \)}{\( 0{.}0015\,\mathrm{m} \)}{}{}}
\LANGpl{
\Question{
Oczekiwano, że drewniane deski zostaną przycięta na tą samą długość. Po ich przycięciu ich wymiary są równe: \( 2{,}00;\ 2{,}02;\ 2{,}05;\ 2{,}02;\ 2{,}08;\ 2{,}11 \) (w metrach). Używamy odchylenia standardowego długości deski, aby opisać  (określić ilościowo) dokładność cięcia. Oblicz odchylenie standardowe długości desek i zaokrągli wynik do czterech miejsc po przecinku.}
{\( 0{,}0382\,\mathrm{m} \)}{\( 0{,}0381\,\mathrm{m} \)}{\( 0{,}0014\,\mathrm{m} \)}{\( 0{,}0015\,\mathrm{m} \)}{}{}}
\LANGcs{
\Question{
Desky měly být nařezané na stejnou délku. Po jejich odřezání a přeměření byly zjištěny následující skutečné délky (v metrech): \( 2{,}00;\  2{,}02;\  2{,}05;\  2{,}02;\  2{,}08;\  2{,}11. \) 
Pro posouzení přesnosti délek bude použita směrodatná odchylka délky desky. Určete směrodatnou odchylku s přesností na čtyři desetinná místa.}
{\( 0{,}0382\,\mathrm{m} \)}{\( 0{,}0381\,\mathrm{m} \)}{\( 0{,}0014\,\mathrm{m} \)}{\( 0{,}0015\,\mathrm{m} \)}{}{}}
\LANGsk{
\Question{
Dosky mali byť narezané na rovnakú dĺžku. Po ich odrezaní a premeraní boli zistené nasledovné skutočné dĺžky (v metroch): \( 2{,}00;\  2{,}02;\  2{,}05;\  2{,}02;\  2{,}08;\  2{,}11. \) 
Pre posúdenie presnosti dĺžok bude použitá smerodajná odchýlka dĺžky dosky. Určte smerodajnú odchýlku s presnosťou na štyri desatinné miesta.}
{\( 0{,}0382\,\mathrm{m} \)}{\( 0{,}0381\,\mathrm{m} \)}{\( 0{,}0014\,\mathrm{m} \)}{\( 0{,}0015\,\mathrm{m} \)}{}{}}
\LANGes{
\Question{
Wooden boards were expected to be cut to the same length. After their cutting the resulting measured lengths are: \( 2{.}00;\ 2{.}02;\ 2{.}05;\ 2{.}02;\ 2{.}08;\ 2{.}11 \) (in meters). We use the standard deviation of board lengths to describe (quantify) the accuracy of cutting. Find the standard deviation of board lengths rounded to four decimal places.}
{\( 0{.}0382\,\mathrm{m} \)}{\( 0{.}0381\,\mathrm{m} \)}{\( 0{.}0014\,\mathrm{m} \)}{\( 0{.}0015\,\mathrm{m} \)}{}{}}
\End

\Data\ID{1003029402}
\Updated{2021-03-11 10:03:17}
\Level{B}
\SubArea{Statistics}
\LANGen{
\Question{
A sample of \( 50 \) pieces of pears was selected randomly from the production of the plant breeding institute. The weights of these pears are recorded in the table.
\[
\begin{array}{|c|c|} \hline \text{ Weight (g) }&\text{ Number of pears } \\\hline 26\text{ -- }30&8 \\\hline31\text{ -- }35&14 \\\hline 36\text{ -- }40&15 \\\hline 41\text{ -- }45&9 \\\hline 46\text{ -- }50&4\\\hline\end{array}\]
Calculate the variance of pear weights and round the variance to the nearest hundredth. 
(To solve the task with help of calculator is recommended.)}
{\( 33{.}81\,\mathrm{g}^2 \)}{\( 5{.}81\,\mathrm{g}^2 \)}{\( 15{.}84\,\mathrm{g}^2 \)}{\( 39{.}84\,\mathrm{g}^2 \)}{}{}}
\LANGpl{
\Question{
\( 50 \) gruszek zostało losowo wybranych przez instytut hodowli roślin. Ich waga została podana w tabeli.
\[
\begin{array}{|c|c|} \hline \text{ Waga (g) }&\text{ Number gruszki } \\\hline 26\text{ -- }30&8 \\\hline31\text{ -- }35&14 \\\hline 36\text{ -- }40&15 \\\hline 41\text{ -- }45&9 \\\hline 46\text{ -- }50&4\\\hline\end{array}\]
Oblicz wariancję wagi gruszek i zaokrągli wynik do dwóch miejsc po przecinku.
(Zaleca się użycie kalkulatora do rozwiązania tego zadania.)}
{\( 33{,}81\,\mathrm{g}^2 \)}{\( 5{,}81\,\mathrm{g}^2 \)}{\( 15{,}84\,\mathrm{g}^2 \)}{\( 39{,}84\,\mathrm{g}^2 \)}{}{}}
\LANGcs{
\Question{
Z úrody hrušek ve šlechtitelském ústavu náhodně vybrali \( 50 \) kusů a zjistili jejich hmotnosti. Výsledky vážení jsou zaznamenané v tabulce:
\[
\begin{array}{|c|c|} \hline \text{ Hmotnost (g) }&\text{ Počet hrušek } \\\hline 26\text{ -- }30&8 \\\hline31\text{ -- }35&14 \\\hline 36\text{ -- }40&15 \\\hline 41\text{ -- }45&9 \\\hline 46\text{ -- }50&4\\\hline\end{array}\]
Jaký je rozptyl hmotností vybraných hrušek? Úlohu řešte pomocí kalkulačky a výsledek zaokrouhlete na dvě desetinná místa.}
{\( 33{,}81\,\mathrm{g}^2 \)}{\( 5{,}81\,\mathrm{g}^2 \)}{\( 15{,}84\,\mathrm{g}^2 \)}{\( 39{,}84\,\mathrm{g}^2 \)}{}{}}
\LANGsk{
\Question{
Z úrody hrušiek v šľachtiteľskom ústave náhodne vybrali \( 50 \) kusov a zistili ich hmotnosti. Výsledky váženia sú zaznamenané v tabuľke:
\[
\begin{array}{|c|c|} \hline \text{ Hmotnosť (g) }&\text{ Počet hrušiek } \\\hline 26\text{ -- }30&8 \\\hline31\text{ -- }35&14 \\\hline 36\text{ -- }40&15 \\\hline 41\text{ -- }45&9 \\\hline 46\text{ -- }50&4\\\hline\end{array}\]
Aký je rozptyl hmotností vybratých hrušiek? Úlohu riešte pomocou kalkulačky a výsledok zaokrúhlite na dve desatinné miesta.}
{\( 33{,}81\,\mathrm{g}^2 \)}{\( 5{,}81\,\mathrm{g}^2 \)}{\( 15{,}84\,\mathrm{g}^2 \)}{\( 39{,}84\,\mathrm{g}^2 \)}{}{}}
\LANGes{
\Question{
En un instituto de crianza han elejido \( 50 \)peras al azar y han anotado sus masas. Los resultados están en la tabla: 
\[
\begin{array}{|c|c|} \hline \text{ Masa (g) }&\text{ Cantidad de peras } \\\hline 26\text{ -- }30&8 \\\hline31\text{ -- }35&14 \\\hline 36\text{ -- }40&15 \\\hline 41\text{ -- }45&9 \\\hline 46\text{ -- }50&4\\\hline\end{array}\]
¿Cómo es la varianza de las masas de las peras? Averigua usando la calculadora y aproxima a dos cifras decimales.}
{\( 33{,}81\,\mathrm{g}^2 \)}{\( 5{,}81\,\mathrm{g}^2 \)}{\( 15{,}84\,\mathrm{g}^2 \)}{\( 39{,}84\,\mathrm{g}^2 \)}{}{}}
\End

\Data\ID{1003134401}
\Updated{2021-03-11 10:03:59}
\Level{B}
\SubArea{Statistics}
\LANGen{
\Question{
We want to compare the performances of two javelin throwers in one competition. Throws of Alex and Martin (in meters) are recorded in the following table. Calculate the coefficient of variation for each set of results and determine, which of the athletes has more balanced performance. I.e. choose the name of the athlete with more balanced performance and with the correct coefficient of variation (\( \% \)) of his throws. The coefficient of variation is rounded to two decimal places.
\[
\begin{array}{|c|c|c|c|c|} \hline \textbf{Alex} & 78.95 & 83.32 & 86.14 & 84.46 \\\hline  \textbf{Martin} & 84.66 & 83.63 & 76.83 & 83.23 \\\hline \end{array} \]}
{Alex: \( 3{.}20\,\% \)}{Alex: \( 27{.}99\,\% \)}{Martin: \( 4{.}52\,\% \)}{Martin: \( 23{.}52\,\% \)}{}{}}
\LANGpl{
\Question{
Chcemy porównać występy dwóch atletów rzucających oszczepem w jednym konkursie. Wyniki  rzutów  Alexa i Martina (w metrach) są zapisane w poniższej tabeli. Oblicz współczynnik zmienności dla każdego zestawu wyników i określ, który sportowiec ma bardziej zrównoważoną wydajność. To znaczy, wybierz imię sportowca o bardziej zrównoważonej wydajności i poprawnym współczynniku zmienności  (\( \% \))  jego rzutów. Współczynnik zmienności jest zaokrąglany do dwóch miejsc po przecinku.
\[
\begin{array}{|c|c|c|c|c|} \hline \textbf{Alex} & 78{,}95 & 83{,}32 & 86{,}14 & 84{,}46 \\\hline  \textbf{Martin} & 84{,}66 & 83{,}63 & 76{,}83 & 83{,}23 \\\hline \end{array} \]}
{Alex: \( 3{,}20\,\% \)}{Alex: \( 27{,}99\,\% \)}{Martin: \( 4{,}52\,\% \)}{Martin: \( 23{,}52\,\% \)}{}{}}
\LANGcs{
\Question{
V tabulce jsou zaznamenané výkony (v metrech) dvou oštěpařů na závodech v atletice. Zjistěte pomocí variačního koeficientu, který závodník podal vyrovnanější výkon. Označte jméno závodníka a variační koeficient jeho výsledků. Variační koeficient je vyjádřený v procentech a zaokrouhlený na dvě desetinná místa.
\[
\begin{array}{|c|c|c|c|c|} \hline \textbf{Alex} & 78{,}95 & 83{,}32 & 86{,}14 & 84{,}46 \\\hline  \textbf{Martin} & 84{,}66 & 83{,}63 & 76{,}83 & 83{,}23 \\\hline \end{array} \]}
{Alex: \( 3{,}20\,\% \)}{Alex: \( 27{,}99\,\% \)}{Martin: \( 4{,}52\,\% \)}{Martin: \( 23{,}52\,\% \)}{}{}}
\LANGsk{
\Question{
V tabuľke sú zaznamenané výkony (v metroch) dvoch oštepárov na pretekoch v atletike. Zistite pomocou variačného koeficientu, ktorý pretekár podal vyrovnanejší výkon. Označte meno pretekára a variačný koeficient jeho výsledkov. Variačný koeficient je vyjadrený v percentách a zaokrúhlený na dve desatinné miesta.
\[
\begin{array}{|c|c|c|c|c|} \hline \textbf{Alex} & 78{,}95 & 83{,}32 & 86{,}14 & 84{,}46 \\\hline  \textbf{Martin} & 84{,}66 & 83{,}63 & 76{,}83 & 83{,}23 \\\hline \end{array} \]}
{Alex: \( 3{,}20\,\% \)}{Alex: \( 27{,}99\,\% \)}{Martin: \( 4{,}52\,\% \)}{Martin: \( 23{,}52\,\% \)}{}{}}
\LANGes{
\Question{
En la tabla hay los resultados (en metros) de lanzamiento de jabalina de dos lanzaderos. Usando el coeficiente de variación averigua cuál de los lanzaderos ha tenido los resultados más uniformes. Elije su nombre y su coeficiente de variación (expresado en porcentaje y aproximado a dos cifras decimales).
\[
\begin{array}{|c|c|c|c|c|} \hline \textbf{Alex} & 78{,}95 & 83{,}32 & 86{,}14 & 84{,}46 \\\hline  \textbf{Martin} & 84{,}66 & 83{,}63 & 76{,}83 & 83{,}23 \\\hline \end{array} \]}
{Alex: \( 3{,}20\,\% \)}{Alex: \( 27{,}99\,\% \)}{Martin: \( 4{,}52\,\% \)}{Martin: \( 23{,}52\,\% \)}{}{}}
\End

\Data\ID{1003134402}
\Updated{2020-07-15 03:07:21}
\Level{B}
\SubArea{Statistics}
\LANGen{
\Question{
There are two groups, A and B, of students in German language class. Each group consists of \( 15 \) students. In the tables, down each column, a student’s ID and grade from the mid-year’s grade report are listed. The students are evaluated according to the grading scale from \( 1 \) to \( 5 \), while \( 1 \) is the best grade and \( 5 \) is the worst one. 

Calculate the coefficient of variation of grades for each group and determine in which group the grades are more balanced. I.e. choose the name of the group with more balanced grades and with the correct coefficient of variation (\( \% \)) of grades. The value of the coefficient of variation is rounded to two decimal places. 

\[ \begin{array}{|c|c|c|c|c|c|c|c|c|}\hline 
\textbf{A -- students} & 1 & 2 & 3 & 4 & 5 & 6 & 7 & 8  \\\hline \textbf{Grade} & 2 & 2 & 2 & 2 & 3 & 2 & 1 & 2  \\\hline \\\hline
\textbf{A -- students} & 9 & 10 & 11 & 12 & 13 & 14 & 15 &  \\\hline \textbf{Grade} &  2 & 1 & 3 & 1 &3 & 2 & 3 & \\\hline
\end{array} \]

\[ \begin{array}{|c|c|c|c|c|c|c|c|c|}\hline 
\textbf{B -- students} & 1 & 2 & 3 & 4 & 5 & 6 & 7 & 8  \\\hline \textbf{Grade} & 2 & 1 & 1 & 2 & 2 & 3 & 1 & 2   \\\hline \\\hline
\textbf{B -- students} & 9 & 10 & 11 & 12 & 13 & 14 & 15 &  \\\hline \textbf{Grade} & 2 & 1 & 2 &1 &1 &1 &1 &  \\\hline
\end{array} \]}
{A: \( 32{.}90\,\% \)}{A: \( 3{.}04\,\% \)}{B: \( 40{.}32\,\% \)}{B: \( 2{.}48\,\% \)}{}{}}
\LANGpl{
\Question{
Na zajęciach z języka niemieckiego są dwie grupy uczniów (A i B). Do każdej grupy uczęszcza  \( 15 \) uczniów. W tabeli, pionowo, znajduje się numer ucznia oraz ocena na półrocze. Uczniowie są oceniani w skali od \( 1 \) do \( 5 \),  \( 1 \) to najlepsza ocena, natomiast  \( 5 \) to najsłabsza ocena.
Oblicz współczynnik zmienności ocen dla każdej grupy i określ, w której grupie oceny są bardziej zrównoważone. To znaczy wybierz nazwę grupy o bardziej zrównoważonych ocenach i poprawnym współczynniku zmienności  (\( \% \)) ocen. Wartość współczynnika zmienności jest  zaokrąglona do dwóch miejsc po przecinku.
\[ \begin{array}{|c|c|c|c|c|c|c|c|c|}\hline 
\textbf{A -- ucznia} & 1 & 2 & 3 & 4 & 5 & 6 & 7 & 8  \\\hline \textbf{Ocena} & 2 & 2 & 2 & 2 & 3 & 2 & 1 & 2  \\\hline \\\hline
\textbf{A -- ucznia} & 9 & 10 & 11 & 12 & 13 & 14 & 15 &  \\\hline \textbf{Ocena} &  2 & 1 & 3 & 1 &3 & 2 & 3 & \\\hline
\end{array} \]

\[ \begin{array}{|c|c|c|c|c|c|c|c|c|}\hline 
\textbf{B -- ucznia} & 1 & 2 & 3 & 4 & 5 & 6 & 7 & 8  \\\hline \textbf{Ocena} & 2 & 1 & 1 & 2 & 2 & 3 & 1 & 2   \\\hline \\\hline
\textbf{B -- ucznia} & 9 & 10 & 11 & 12 & 13 & 14 & 15 &  \\\hline \textbf{Ocena} & 2 & 1 & 2 &1 &1 &1 &1 &  \\\hline
\end{array} \]}
{A: \( 32{,}90\,\% \)}{A: \( 3{,}04\,\% \)}{B: \( 40{,}32\,\% \)}{B: \( 2{,}48\,\% \)}{}{}}
\LANGcs{
\Question{
Žáci jedné třídy jsou na hodině německého jazyka rozdělení na skupiny A a B po \( 15 \) žácích. V tabulkách jsou uvedené jejich známky za půl roku (žáci jsou hodnocení na stupnici \( 1 \) - \( 5 \), kde \( 1 \) je nejlepší hodnocení a \( 5 \) nejhorší hodnocení). Zjistěte pomocí variačního koeficientu, která skupina dosáhla v německém jazyku vyrovnanějších výsledků. Označte číslo skupiny a variační koeficient známek studentů teto skupiny. Variační koeficient je vyjádřený v procentech a zaokrouhlený na dvě desetinná místa. Použijte na výpočty statistický režim kalkulačky.

\[ \begin{array}{|c|c|c|c|c|c|c|c|c|}\hline 
\textbf{A -- žáci} & 1 & 2 & 3 & 4 & 5 & 6 & 7 & 8  \\\hline \textbf{Známka} & 2 & 2 & 2 & 2 & 3 & 2 & 1 & 2  \\\hline \\\hline
\textbf{A -- žáci} & 9 & 10 & 11 & 12 & 13 & 14 & 15 &  \\\hline \textbf{Známka} &  2 & 1 & 3 & 1 &3 & 2 & 3 & \\\hline
\end{array} \]

\[ \begin{array}{|c|c|c|c|c|c|c|c|c|}\hline 
\textbf{B -- žáci} & 1 & 2 & 3 & 4 & 5 & 6 & 7 & 8  \\\hline \textbf{Známka} & 2 & 1 & 1 & 2 & 2 & 3 & 1 & 2   \\\hline \\\hline
\textbf{B -- žáci} & 9 & 10 & 11 & 12 & 13 & 14 & 15 &  \\\hline \textbf{Známka} & 2 & 1 & 2 &1 &1 &1 &1 &  \\\hline
\end{array} \]}
{A: \( 32{,}90\,\% \)}{A: \( 3{,}04\,\% \)}{B: \( 40{,}32\,\% \)}{B: \( 2{,}48\,\% \)}{}{}}
\LANGsk{
\Question{
Žiaci jednej triedy sú na hodiny nemeckého jazyka rozdelení na skupiny A a B po \( 15 \) žiakov. V tabuľkách sú uvedené ich známky na pol roku (žiaci sú hodnotení na stupnici \( 1 \) - \( 5 \), kde \( 1 \) je najlepšie hodnotenie a \( 5 \) najhoršie hodnotenie). Zistite pomocou variačného koeficientu, ktorá skupina dosiahla v nemeckom jazyku vyrovnanejšie výsledky. Označte číslo skupiny a variačný koeficient známok študentov tejto skupiny. Variačný koeficient je vyjadrený v percentách a zaokrúhlený na dve desatinné miesta. Použite na výpočty štatistický režim kalkulačky.
\[ \begin{array}{|c|c|c|c|c|c|c|c|c|}\hline 
\textbf{A -- žiaci} & 1 & 2 & 3 & 4 & 5 & 6 & 7 & 8  \\\hline \textbf{Známka} & 2 & 2 & 2 & 2 & 3 & 2 & 1 & 2  \\\hline \\\hline
\textbf{A -- žiaci} & 9 & 10 & 11 & 12 & 13 & 14 & 15 &  \\\hline \textbf{Známka} &  2 & 1 & 3 & 1 &3 & 2 & 3 & \\\hline
\end{array} \]

\[ \begin{array}{|c|c|c|c|c|c|c|c|c|}\hline 
\textbf{B -- žiaci} & 1 & 2 & 3 & 4 & 5 & 6 & 7 & 8  \\\hline \textbf{Známka} & 2 & 1 & 1 & 2 & 2 & 3 & 1 & 2   \\\hline \\\hline
\textbf{B -- žiaci} & 9 & 10 & 11 & 12 & 13 & 14 & 15 &  \\\hline \textbf{Známka} & 2 & 1 & 2 &1 &1 &1 &1 &  \\\hline
\end{array} \]}
{A: \( 32{,}90\,\% \)}{A: \( 3{,}04\,\% \)}{B: \( 40{,}32\,\% \)}{B: \( 2{,}48\,\% \)}{}{}}
\LANGes{
\Question{
There are two groups, A and B, of students in German language class. Each group consists of \( 15 \) students. In the tables, down each column, a student’s ID and grade from the mid-year’s grade report are listed. The students are evaluated according to the grading scale from \( 1 \) to \( 5 \), while \( 1 \) is the best grade and \( 5 \) is the worst one. 

Calculate the coefficient of variation of grades for each group and determine in which group the grades are more balanced. I.e. choose the name of the group with more balanced grades and with the correct coefficient of variation (\( \% \)) of grades. The value of the coefficient of variation is rounded to two decimal places. 

\[ \begin{array}{|c|c|c|c|c|c|c|c|c|}\hline 
\textbf{A -- students} & 1 & 2 & 3 & 4 & 5 & 6 & 7 & 8  \\\hline \textbf{Grade} & 2 & 2 & 2 & 2 & 3 & 2 & 1 & 2  \\\hline \\\hline
\textbf{A -- students} & 9 & 10 & 11 & 12 & 13 & 14 & 15 &  \\\hline \textbf{Grade} &  2 & 1 & 3 & 1 &3 & 2 & 3 & \\\hline
\end{array} \]

\[ \begin{array}{|c|c|c|c|c|c|c|c|c|}\hline 
\textbf{B -- students} & 1 & 2 & 3 & 4 & 5 & 6 & 7 & 8  \\\hline \textbf{Grade} & 2 & 1 & 1 & 2 & 2 & 3 & 1 & 2   \\\hline \\\hline
\textbf{B -- students} & 9 & 10 & 11 & 12 & 13 & 14 & 15 &  \\\hline \textbf{Grade} & 2 & 1 & 2 &1 &1 &1 &1 &  \\\hline
\end{array} \]}
{A: \( 32{.}90\,\% \)}{A: \( 3{.}04\,\% \)}{B: \( 40{.}32\,\% \)}{B: \( 2{.}48\,\% \)}{}{}}
\End

\Data\ID{1003134403}
\Updated{2021-03-11 09:03:57}
\Level{B}
\SubArea{Statistics}
\LANGen{
\Question{
The average age of town citizens decreased by \( 19\,\% \) due to a satellite town construction. The variance of the age has increased by \( 21\,\% \).  Complete the correct statement. The coefficient of variation .... (Note: The results are rounded to two decimal places.)}
{increased by \( 35{.}80\,\% \).}{increased by \( 49{.}38\,\% \).}{decreased by \( 33{.}06\,\% \).}{decreased by \( 26{.}36\,\% \).}{}{}}
\LANGpl{
\Question{
Średni wiek mieszkańców miasta zmniejszył się o  \( 19\,\% \) z powodu budowy miasta satelitarnego. Wariancja wieku wzrosła o  \( 21\,\% \).  Uzupełnij poprawnie zdanie. Współczynnik zmienności .... (Uwaga: wyniki są zaokrąglone do dwóch miejsc po przecinku.)}
{zwiększył się o \( 35{,}80\,\% \).}{zwiększył się o  \( 49{,}38\,\% \).}{zmniejszył się o  \( 33{,}06\,\% \).}{zmniejszył się o \( 26{,}36\,\% \).}{}{}}
\LANGcs{
\Question{
V důsledku výstavby satelitního městečka poklesl průměrný věk obyvatelů obce o \( 19\,\% \), rozptyl věku vzrostl o  \( 21\,\% \).  Jak se změnil variační koeficient? Výsledky jsou zaokrouhlené na dvě desetinná místa.}
{Vzrostl o \( 35{,}80\,\% \).}{Vzrostl o \( 49{,}38\,\% \).}{Poklesl o \( 33{,}06\,\% \).}{Poklesl o \( 26{,}36\,\% \).}{}{}}
\LANGsk{
\Question{
V dôsledku výstavby satelitného mestečka poklesol priemerný vek obyvateľov obce o \( 19\,\% \), rozptyl veku vzrástol o  \( 21\,\% \).  Ako sa zmenil variačný koeficient? Výsledky sú zaokrúhlené na dve desatinné miesta.}
{Vzrástol o \( 35{,}80\,\% \).}{Vzrástol o \( 49{,}38\,\% \).}{Poklesol o \( 33{,}06\,\% \).}{Poklesol o \( 26{,}36\,\% \).}{}{}}
\LANGes{
\Question{
La edad media de los ciudadanos ha disminuido por  \( 19\,\% \), y la varianza de la edad ha aumentado por \( 21\,\% \) por la construcción de la ciudad satélite. ¿Cómo ha cambiado el coeficiente de variación? Aproxima los resultados a dos cifras decimales.}
{Ha aumentado por \( 35{,}80\,\% \).}{Ha aumentado por \( 49{,}38\,\% \).}{Ha disminuido por \( 33{,}06\,\% \).}{Ha disminuido por \( 26{,}36\,\% \).}{}{}}
\End

\Data\ID{1003134407}
\Updated{2021-03-11 09:03:39}
\Level{B}
\SubArea{Statistics}
\LANGen{
\Question{
In the tables below the absent hours in lessons of \( 16 \) boys and \( 14 \) girls from one class for half of a year are listed.  Use the variance of the number of absent hours to find out the students of which gender had the absence more balanced. I.e. choose the group with more balanced absence and the correct variance of the number of absent hours. The variance is rounded to two decimal places.
\[ \begin{array}{|c|c|c|c|c|c|c|c|} \hline 
\text{Girl's ID} & 1 & 2 & 3 & 4 & 5 & 6 & 7  \\\hline \text{Absent hours} & 27 & 61 & 38 & 61 & 17 & 39 & 61 \\\hline \\\hline
\text{Girl's ID} & 8 & 9 & 10 & 11 & 12 & 13 & 14 \\\hline \text{Absent hours} & 25 & 21 & 52 & 16 & 34 & 9 & 25 \\\hline \end{array} \]
\[  \begin{array}{|c|c|c|c|c|c|c|c|c|c|c|} \hline 
\text{Boy's ID} & 1 & 2 & 3 & 4 & 5 & 6 & 7 & 8  \\\hline \text{Absent hours} & 67 & 56 & 26 & 36 & 27 & 55 & 17 & 34  \\\hline \\\hline
\text{Boy's ID} & 9 & 10 & 11 & 12 & 13 & 14 & 15 & 16 \\\hline \text{Absent hours} & 54 & 46 & 13 & 48 & 21 & 49 & 18 & 14 \\\hline
 \end{array} \]}
{boys: \( \sigma^2= 285{.}34\,\mathrm{lessons}^2 \)}{girls: \( \sigma^2= 297{.}35\,\mathrm{lessons}^2 \)}{boys: \( \sigma^2= 16{.}89\,\mathrm{lessons} \)}{girls: \( \sigma^2= 17{.}24\,\mathrm{lessons} \)}{}{}}
\LANGpl{
\Question{
W tabelach poniżej podano nieobecności na lekcjach \( 16 \) chłopców oraz  \( 14 \) dziewczyn z jednej klasy w ciągu jednego półrocza. Użyj wariancji liczby nieobecności godzin, aby dowiedzieć się która z płci ma bardziej zrównoważoną nieobecność tz.  wybierz grupę o bardziej zrównoważonej nieobecności i poprawnej wariancji liczby nieobecnych godzin. Wariancja jest zaokrąglona do dwóch miejsc po przecinku.
\[ \begin{array}{|c|c|c|c|c|c|c|c|} \hline 
\text{ID dziewczyny} & 1 & 2 & 3 & 4 & 5 & 6 & 7  \\\hline \text{Liczba nieobecnych godzin} & 27 & 61 & 38 & 61 & 17 & 39 & 61 \\\hline \\\hline
\text{ID dziewczyny} & 8 & 9 & 10 & 11 & 12 & 13 & 14 \\\hline \text{Liczba nieobecnych godzin} & 25 & 21 & 52 & 16 & 34 & 9 & 25 \\\hline \end{array} \]
\[  \begin{array}{|c|c|c|c|c|c|c|c|c|c|c|} \hline 
\text{ID chłopca} & 1 & 2 & 3 & 4 & 5 & 6 & 7 & 8  \\\hline \text{Liczba nieobecnych godzin} & 67 & 56 & 26 & 36 & 27 & 55 & 17 & 34  \\\hline \\\hline
\text{ID chłopca} & 9 & 10 & 11 & 12 & 13 & 14 & 15 & 16 \\\hline \text{Liczba nieobecnych godzin} & 54 & 46 & 13 & 48 & 21 & 49 & 18 & 14 \\\hline
 \end{array} \]}
{chłopcy: \( \sigma^2= 285{,}34\,\mathrm{lekcji}^2 \)}{dziewczyny: \( \sigma^2= 297{,}35\,\mathrm{lekcji}^2 \)}{chłopcy: \( \sigma^2= 16{,}89\,\mathrm{lekcji} \)}{dziewczyny: \( \sigma^2= 17{,}24\,\mathrm{lekcji} \)}{}{}}
\LANGcs{
\Question{
V tabulkách jsou uvedené zameškané hodiny chlapců a děvčat jedné třídy za jeden půlrok. Zjistěte pomocí rozptylu  \( \sigma^2 \), která skupina měla rovnoměrnější absenci. Označte tuto skupinu a její rozptyl zameškaných hodin zaokrouhlený na dvě desetinná místa. Pro výpočty použijte statistický režim kalkulačky. Výsledky zaokrouhlete na dvě desetinná místa.
\[ \begin{array}{|c|c|c|c|c|c|c|c|} \hline 
\text{ID dívky} & 1 & 2 & 3 & 4 & 5 & 6 & 7  \\\hline \text{Počet hodin} & 27 & 61 & 38 & 61 & 17 & 39 & 61 \\\hline \\\hline
\text{ID dívky} & 8 & 9 & 10 & 11 & 12 & 13 & 14 \\\hline \text{Počet hodin} & 25 & 21 & 52 & 16 & 34 & 9 & 25 \\\hline \end{array} \]
\[  \begin{array}{|c|c|c|c|c|c|c|c|c|c|c|} \hline 
\text{ID chlapce} & 1 & 2 & 3 & 4 & 5 & 6 & 7 & 8  \\\hline \text{Počet hodín} & 67 & 56 & 26 & 36 & 27 & 55 & 17 & 34  \\\hline \\\hline
\text{ID chlapce} & 9 & 10 & 11 & 12 & 13 & 14 & 15 & 16 \\\hline \text{Počet hodín} & 54 & 46 & 13 & 48 & 21 & 49 & 18 & 14 \\\hline
 \end{array} \]}
{chlapci: \( \sigma^2= 285{,}34\,\text{hodin}^2 \)}{děvčata: \( \sigma^2= 297{,}35\,\text{hodin}^2 \)}{chlapci: \( \sigma^2= 16{,}89\,\text{hodin} \)}{děvčata: \( \sigma^2= 17{,}24\,\text{hodin} \)}{}{}}
\LANGsk{
\Question{
V tabuľkách sú uvedené vymeškané hodiny chlapcov a dievčat jednej triedy za jeden pol rok. Zistite pomocou rozptylu  \( \sigma^2 \) , ktorá skupina mala rovnomernejšiu absenciu. Označte túto skupinu a jej rozptyl vymeškaných hodín zaokrúhlený na dve desatinné miesta. Na výpočty použite štatistický režim kalkulačky. Výsledky zaokrúhlite na dve desatinné miesta.
\[ \begin{array}{|c|c|c|c|c|c|c|c|} \hline 
\text{ID dievča} & 1 & 2 & 3 & 4 & 5 & 6 & 7  \\\hline \text{počet hodín} & 27 & 61 & 38 & 61 & 17 & 39 & 61 \\\hline \\\hline
\text{ID dievča} & 8 & 9 & 10 & 11 & 12 & 13 & 14 \\\hline \text{počet hodín} & 25 & 21 & 52 & 16 & 34 & 9 & 25 \\\hline \end{array} \]
\[  \begin{array}{|c|c|c|c|c|c|c|c|c|c|c|} \hline 
\text{ID chlapca} & 1 & 2 & 3 & 4 & 5 & 6 & 7 & 8  \\\hline \text{počet hodín} & 67 & 56 & 26 & 36 & 27 & 55 & 17 & 34  \\\hline \\\hline
\text{ID chlapca} & 9 & 10 & 11 & 12 & 13 & 14 & 15 & 16 \\\hline \text{počet hodín} & 54 & 46 & 13 & 48 & 21 & 49 & 18 & 14 \\\hline
 \end{array} \]}
{chlapci: \( \sigma^2= 285{,}34\,\text{hodín}^2 \)}{dievčatá: \( \sigma^2= 297{,}35\,\text{hodín}^2 \)}{chlapci: \( \sigma^2= 16{,}89\,\text{hodín} \)}{dievčatá: \( \sigma^2= 17{,}24\,\text{hodín} \)}{}{}}
\LANGes{
\Question{
En las tablas hay las horas de ausencia de las chicas y los chicos de una clase durante un año. Usando la varianza \( \sigma^2 \), averigua cuál de los grupos ha tenido ausencia máa uniforme. Elije el grupo y su varianza aproximada a dos cifras decimales. 
\[ \begin{array}{|c|c|c|c|c|c|c|c|} \hline 
\text{ID de chica} & 1 & 2 & 3 & 4 & 5 & 6 & 7  \\\hline \text{Cantidad de horas} & 27 & 61 & 38 & 61 & 17 & 39 & 61 \\\hline \\\hline
\text{ID de chica} & 8 & 9 & 10 & 11 & 12 & 13 & 14 \\\hline \text{Cantidad de horas} & 25 & 21 & 52 & 16 & 34 & 9 & 25 \\\hline \end{array} \]
\[  \begin{array}{|c|c|c|c|c|c|c|c|c|c|c|} \hline 
\text{ID de chico} & 1 & 2 & 3 & 4 & 5 & 6 & 7 & 8  \\\hline \text{Cantidad de horas} & 67 & 56 & 26 & 36 & 27 & 55 & 17 & 34  \\\hline \\\hline
\text{ID de chico} & 9 & 10 & 11 & 12 & 13 & 14 & 15 & 16 \\\hline \text{Cantidad de horas} & 54 & 46 & 13 & 48 & 21 & 49 & 18 & 14 \\\hline
 \end{array} \]}
{chicos: \( \sigma^2= 285{,}34\,\text{horas}^2 \)}{chicas: \( \sigma^2= 297{,}35\,\text{horas}^2 \)}{chicos: \( \sigma^2= 16{,}89\,\text{horas} \)}{chicas: \( \sigma^2= 17{,}24\,\text{horas} \)}{}{}}
\End

\Data\ID{1103134404}
\Updated{2021-05-04 12:05:37}
\Level{B}
\SubArea{Statistics}
\IMG{1103134404.pdf}
\LANGen{
\Question{
The diagram shows the values of four continuous variables (A, B, C and D). Determine which one of these variables has the least coefficient of variation.}
{D}{C}{B}{A}{}{}}
\LANGpl{
\Question{
Wykres pokazuje wartości czterech zmiennych ciągłych (A, B, C and D). Wskaż, która z tych zmiennych ma najmniejszy współczynnik zmienności.}
{D}{C}{B}{A}{}{}}
\LANGcs{
\Question{
Na obrázku jsou prezentované hodnoty čtyř spojitých znaků (A, B, C a D). Určete, který z těchto znaků má nejmenší variační koeficient.}
{D}{C}{B}{A}{}{}}
\LANGsk{
\Question{
Na obrázku sú prezentované hodnoty štyroch spojitých znakov (A, B, C a D). Určite, ktorý z týchto znakov má najmenší variačný koeficient.}
{D}{C}{B}{A}{}{}}
\LANGes{
\Question{
En el dibujo están representadas las medidas de cuatro variables contínuas (A, B, C y D). Uverigua cuál de las variables tiene el menor coeficiente de variación.}
{D}{C}{B}{A}{}{}}
\End

\Data\ID{1103134405}
\Updated{2020-07-15 03:07:21}
\Level{B}
\SubArea{Statistics}
\IMG{1103134405.pdf}
\LANGen{
\Question{
The students are evaluated according to the grading scale from \( 1 \) to \( 5 \), while \( 1 \) is the best grade and \( 5 \) is the worst one. In the pictures, there are visualizations of relative frequencies of grades in Math, which students of two groups (A and B) had in their year-class-reports. Calculate the variance of grades for each group of students and determine, in which group student performances of Math knowledge are more balanced. I.e. from the offered choices choose the group, which has more balanced grades and with the correct variance of grades. The variance is rounded to two decimal places.}
{A: \( 0{.}81 \)}{B: \( 0{.}84 \)}{A: \( 0{.}90 \)}{B: \( 0{.}92 \)}{}{}}
\LANGpl{
\Question{
Uczniowie są oceniani według skali ocen od \( 1 \) do  \( 5 \),  \( 1 \) to najlepsza ocena  \( 5 \) to najsłabsza ocena.  Na zdjęciach są wizualizacje względnych częstotliwości ocen z matematyki, które uczniowie z dwóch grup (A i  B) otrzymali na koniec roku. Oblicz wariancję ocen dla każdej grupy uczniów i określ, w której grupie wyniki uczniów z matematyki są bardziej zrównoważone tz. z dostępnych opcji wybierz grupę, która ma bardziej zrównoważone oceny i poprawną wariancję ocen. Wariancja jest zaokrąglona do dwóch miejsc po przecinku.}
{A: \( 0{,}81 \)}{B: \( 0{,}84 \)}{A: \( 0{,}90 \)}{B: \( 0{,}92 \)}{}{}}
\LANGcs{
\Question{
Žáci jsou hodnocení na stupnici \( 1 \) - \( 5 \), kde je \( 1 \) nejlepší hodnocení a \( 5 \) nejhorší hodnocení. Na obrázcích jsou graficky zobrazené relativní četnosti známek z matematiky, které na vysvědčení měli žáci ve dvou třídách (A a B) v jednom ročníku. Určete, ve které třídě dosáhli žáci v matematice vyrovnanějších vědomostí. Označte tuto třídu a rozptyl známek jejích studentů. Rozptyl zaokrouhlujte na dvě desetinná místa.
{Poznámka: Na obrázku "Grade" znamená "Známka".}}
{A: \( 0{,}81 \)}{B: \( 0{,}84 \)}{A: \( 0{,}90 \)}{B: \( 0{,}92 \)}{}{}}
\LANGsk{
\Question{
Žiaci sú hodnotení na stupnici \( 1 \) - \( 5 \), kde je \( 1 \) najlepšie hodnotenie a \( 5 \) najhoršie hodnotenie. Na obrázkoch sú graficky zobrazené relatívne početnosti známok z matematiky, ktoré na vysvedčení mali žiaci v dvoch triedach (A a B) v jednom ročníku. Určte, v ktorej triede dosiahli žiaci v matematike vyrovnanejšie vedomosti. Označte tuto triedu a rozptyl známok ich študentov. Rozptyl zaokrúhľujte na dve desatinné miesta.
{Poznámka: Na obrázku "Grade" znamená Známka.}}
{A: \( 0{,}81 \)}{B: \( 0{,}84 \)}{A: \( 0{,}90 \)}{B: \( 0{,}92 \)}{}{}}
\LANGes{
\Question{
The students are evaluated according to the grading scale from \( 1 \) to \( 5 \), while \( 1 \) is the best grade and \( 5 \) is the worst one. In the pictures, there are visualizations of relative frequencies of grades in Math, which students of two groups (A and B) had in their year-class-reports. Calculate the variance of grades for each group of students and determine, in which group student performances of Math knowledge are more balanced. I.e. from the offered choices choose the group, which has more balanced grades and with the correct variance of grades. The variance is rounded to two decimal places.}
{A: \( 0{.}81 \)}{B: \( 0{.}84 \)}{A: \( 0{.}90 \)}{B: \( 0{.}92 \)}{}{}}
\End

\Data\ID{1103134406}
\Updated{2021-05-04 12:05:28}
\Level{B}
\SubArea{Statistics}
\IMG{1103134406.pdf}
\LANGen{
\Question{
The diagram shows the values of four continuous variables (A, B, C and D). Determine which one of these variables has the largest variance.}
{C}{B}{A}{D}{}{}}
\LANGpl{
\Question{
Wykres pokazuje wartości czterech zmiennych ciągłych (A, B, C i D). Określ, która z tych zmiennych ma największą wariancję.}
{C}{B}{A}{D}{}{}}
\LANGcs{
\Question{
Na obrázku jsou prezentované hodnoty čtyř spojitých znaků (A, B, C a D). Určete, který z těchto znaků má největší rozptyl.}
{C}{B}{A}{D}{}{}}
\LANGsk{
\Question{
Na obrázku sú prezentované hodnoty štyroch spojitých znakov (A, B, C a D). Určite, ktorý z týchto znakov má najväčší rozptyl.}
{C}{B}{A}{D}{}{}}
\LANGes{
\Question{
En el dibujo están repreentadas las medidas de cuatro variables contínuas (A, B, C y D). Uverigua cuál de las variables tene la mayor dispersión.}
{C}{B}{A}{D}{}{}}
\End

\Data\ID{9000153301}
\Updated{2020-07-15 03:07:37}
\Level{B}
\SubArea{Statistics}
\LANGen{
\Question{
A student performed repeated measurements of a length (in meters) and evaluated
the main statistical characteristics: mean, standard deviation, variance and the
coefficient of variation. Which of these characteristics has it's unit square meter?}
{variance}{standard deviation}{mean}{coefficient of variation}{}{}}
\LANGpl{
\Question{
Student przeprowadził wielokrotne pomiary długości (w metrach) i ocenił
główne cechy statystyczne: średnią, odchylenie standardowe, wariancję i
współczynnik zmienności. Która z tych cech ma jednostkę metra kwadratowego?}
{wariancja}{odchylenie standardowe}{średnia}{współczynnik zmienności}{}{}}
\LANGcs{
\Question{
Student opakovaně měřil délku tělesa (v metrech). Naměřené hodnoty měl
statisticky zpracovat a vypočítat aritmetický průměr, směrodatnou odchylku,
rozptyl a variační koeficient měření. Která z těchto charakteristik má
jednotku \( \mathrm{m}^{2}\)?}
{rozptyl}{směrodatná odchylka}{aritmetický průměr}{variační koeficient}{}{}}
\LANGsk{
\Question{
Študent opakovane meral dĺžku telesa (v metroch). Namerané hodnoty mal štatisticky spracovať a vypočítať aritmetický priemer, smerodajnú odchýlku, rozptyl a variačný koeficient merania. Ktorá z týchto charakteristík má jednotku štvorcový meter?}
{rozptyl}{smerodajná odchýlka}{aritmetický priemer}{variačný koeficient}{}{}}
\LANGes{
\Question{
A student performed repeated measurements of a length (in meters) and evaluated
the main statistical characteristics: mean, standard deviation, variance and the
coefficient of variation. Which of these characteristics has it's unit square meter?}
{variance}{standard deviation}{mean}{coefficient of variation}{}{}}
\End

\Data\ID{9000153304}
\Updated{2020-07-15 03:07:37}
\Level{B}
\SubArea{Statistics}
\LANGen{
\Question{
Two students performed repeated measurements of the same length. After
evaluating statistical characteristics they got equal means of the length. Which of the
following statement is true? (Remark: The accuracy is determined by the coefficient of variation.)}
{There is not enough information to compare the accuracy of both measurements.}{One of the students measured with a better accuracy than the other.}{Both students did their measurements with the same accuracy.}{}{}{}}
\LANGpl{
\Question{
Dwóch uczniów przeprowadziło wielokrotne pomiary tej samej długości, po  ocenie
głównych  cech statystycznych otrzymali te same średnie długości. Oznacz zdanie prawdziwe. 
 (Uwaga: Dokładność jest określona przez współczynnik zmienności.)}
{Brak wystarczających informacji, aby porównać dokładność obu pomiarów.}{Jeden z uczniów dokonał bardziej dokładnego pomiaru.}{Obu uczniów dokonało pomiaru z taką samą dokładnością.}{}{}{}}
\LANGcs{
\Question{
Dva studenti měřili délku stejného tělesa. Při zpracování naměřených
hodnot zjistili, že mají naprosto stejné aritmetické průměry. Vyberte
pravdivé tvrzení o přesnosti měření obou studentů. (Poznámka: Za míru přesnosti měření považujte jeho relativní chybu
vyjádřenou variačním koeficientem.)}
{Z daných informací nemůžeme jednoznačně rozhodnout, jestli studenti
měřili se stejnou přesností.}{Jeden ze studentů musel určitě měřit s vyšší přesností.}{Přesnost obou studentů byla stejná.}{}{}{}}
\LANGsk{
\Question{
Dvaja študenti merali dĺžku rovnakého telesa. Pri spracovaní nameraných hodnôt zistili, že majú úplne rovnaké aritmetické priemery. Vyberte pravdivé tvrdenie o presnosti merania oboch študentov. (Poznámka: Za mieru presnosti merania považujte jeho relatívnu chybu vyjadrenú variačným koeficientom.)}
{Z daných informácií nemôžeme jednoznačne rozhodnúť, či študenti merali s rovnakou presnosťou.}{Jeden zo študentov meral s vyššou presnosťou.}{Obaja študenti merali s rovnakou presnosťou.}{}{}{}}
\LANGes{
\Question{
Two students performed repeated measurements of the same length. After
evaluating statistical characteristics they got equal means of the length. Which of the
following statement is true? (Remark: The accuracy is determined by the coefficient of variation.)}
{There is not enough information to compare the accuracy of both measurements.}{One of the students measured with a better accuracy than the other.}{Both students did their measurements with the same accuracy.}{}{}{}}
\End

\Data\ID{9000153305}
\Updated{2020-07-15 03:07:37}
\Level{B}
\SubArea{Statistics}
\LANGen{
\Question{
Two students performed repeated measurements of the same length. After
evaluating statistical characteristics they got equal standard deviations. Which of the
following statement is true? (Remark: The accuracy is determined by the coefficient of variation.)}
{There is not enough information to compare the accuracy of both measurements.}{One of the students measured with a better accuracy than the other.}{Both students did their measurements with the same accuracy.}{}{}{}}
\LANGpl{
\Question{
Dwóch uczniów przeprowadziło wielokrotne pomiary tej samej długości, po ocenie głównych cech statystycznych otrzymali to samo odchylenie standardowe. Oznacz zdanie prawdziwe.  (Uwaga: Dokładność jest określona przez współczynnik zmienności.)}
{Brak wystarczających informacji, aby porównać dokładność obu pomiarów.}{Jeden z uczniów dokonał bardziej dokładnego pomiaru.}{Obu uczniów dokonało pomiaru z taką samą dokładnością.}{}{}{}}
\LANGcs{
\Question{
Dva studenti měřili délku stejného tělesa. Při zpracování naměřených
hodnot zjistili, že mají naprosto stejné směrodatné odchylky. Vyberte
pravdivé tvrzení o přesnosti měření obou studentů. (Poznámka: Za míru přesnosti měření považujte jeho relativní chybu
vyjádřenou variačním koeficientem.)}
{Z daných informací nemůžeme jednoznačně rozhodnout, jestli studenti
měřili se stejnou přesností.}{Jeden ze studentů musel určitě měřit s vyšší přesností.}{Přesnost obou studentů byla stejná.}{}{}{}}
\LANGsk{
\Question{
Dvaja študenti merali dĺžku rovnakého telesa. Pri spracovaní nameraných hodnôt zistili, že majú úplne rovnaké smerodajné odchýlky. Vyberte pravdivé tvrdenie o presnosti merania oboch študentov. (Poznámka: Za mieru presnosti merania považujte jeho relatívnu chybu vyjadrenú variačným koeficientom.)}
{Z daných informácií nemôžeme jednoznačne rozhodnúť, či študenti merali s rovnakou presnosťou.}{Jeden zo študentov meral s vyššou presnosťou.}{Obaja študenti merali s rovnakou presnosťou.}{}{}{}}
\LANGes{
\Question{
Two students performed repeated measurements of the same length. After
evaluating statistical characteristics they got equal standard deviations. Which of the
following statement is true? (Remark: The accuracy is determined by the coefficient of variation.)}
{There is not enough information to compare the accuracy of both measurements.}{One of the students measured with a better accuracy than the other.}{Both students did their measurements with the same accuracy.}{}{}{}}
\End

\Data\ID{9000153306}
\Updated{2020-07-15 03:07:37}
\Level{B}
\SubArea{Statistics}
\LANGen{
\Question{
Two students performed repeated measurements of the same length. After
evaluating statistical characteristics they got equal means and also equal standard
deviations. Which of the following statement is true? (Remark: The accuracy is determined by the coefficient of variation.)}
{Both students did their measurements with the same accuracy.}{There is not enough information to compare the accuracy of both measurements.}{One of the students measured with a better accuracy than the other.}{The stated question is irrelevant, two different statistical files cannot share their
means and their standard deviations.}{}{}}
\LANGpl{
\Question{
Dwóch uczniów przeprowadziło wielokrotne pomiary tej samej długości, po ocenie  cech statystycznych otrzymali te same średnie i to samo odchylenie standardowe. Oznacz zdanie prawdziwe. (Uwaga: Dokładność jest określona przez współczynnik zmienności.)}
{Obu uczniów dokonało pomiaru z taką samą dokładnością.}{Brak wystarczających informacji, aby porównać dokładność obu pomiarów.}{Jeden z uczniów dokonał bardziej dokładnego pomiaru.}{Dwa różne układy  statystyczne nie mogą mieć tych samych średnich i odchyleń standardowych.}{}{}}
\LANGcs{
\Question{
Dva studenti měřili délku stejného tělesa. Jejich statistické soubory nebyly
totožné, přesto při zpracování naměřených hodnot zjistili, že mají
naprosto stejné aritmetické průměry i směrodatné odchylky. Vyberte
pravdivé tvrzení o přesnosti měření obou studentů. (Poznámka: Za míru přesnosti měření považujte jeho relativní chybu
vyjádřenou variačním koeficientem.)}
{Přesnost měření obou studentů byla stejná.}{Z daných informací nemůžeme jednoznačně rozhodnout, jestli studenti
měřili se stejnou přesností.}{Jeden ze studentů musel určitě měřit přesněji.}{Otázkou přesností měření se nemá smysl zabývat, neboť nejsou-li
statistické soubory totožné, nemohou mít stejné aritmetické průměry i
směrodatné odchylky.}{}{}}
\LANGsk{
\Question{
Dvaja študenti merali dĺžku rovnakého telesa. Ich štatistické súbory neboli totožné, napriek tomu pri spracovaní nameraných hodnôt zistili, že majú úplne rovnaké aritmetické priemery aj smerodajné odchýlky. Vyberte pravdivé tvrdenie o presnosti merania oboch študentov. 
(Poznámka: Za mieru presnosti merania považujte jeho relatívnu chybu vyjadrenú variačným koeficientom.)}
{Obaja študenti merali s rovnakou presnosťou.}{Z daných informácií nemôžeme jednoznačne rozhodnúť, či študenti merali s rovnakou presnosťou.}{Jeden zo študentov meral s vyššou presnosťou.}{Otázkou presnosti merania nemá zmysel sa zaoberať, pretože ak nie sú štatistické súbory totožné, nemôžu mať rovnaké aritmetické priemery aj smerodajné odchýlky.}{}{}}
\LANGes{
\Question{
Two students performed repeated measurements of the same length. After
evaluating statistical characteristics they got equal means and also equal standard
deviations. Which of the following statement is true? (Remark: The accuracy is determined by the coefficient of variation.)}
{Both students did their measurements with the same accuracy.}{There is not enough information to compare the accuracy of both measurements.}{One of the students measured with a better accuracy than the other.}{The stated question is irrelevant, two different statistical files cannot share their
means and their standard deviations.}{}{}}
\End

\Data\ID{9000153308}
\Updated{2020-07-15 03:07:37}
\Level{B}
\SubArea{Statistics}
\LANGen{
\Question{
Statistical file contains repeating measurements of a body mass in kilograms. Find
the change in the coefficient of variation if we convert all data in the file to grams.}
{The coefficient of variation does not change.}{The coefficient of variation becomes bigger.}{The coefficient of variation becomes smaller.}{}{}{}}
\LANGpl{
\Question{
Zestaw statystyczny zawiera dane dotyczące wielokrotnego pomiaru wagi opakowań mąki w kilogramach. Jak zmieni się współczynnik zmienności, jeśli masa opakowania będzie  podana w gramach?}
{Współczynnik zmienności nie zmieni się.}{Współczynnik zmienności wzrośnie.}{Współczynnik zmienności zmniejszy się.}{}{}{}}
\LANGcs{
\Question{
Statistický soubor obsahuje údaje o opakovaném měření hmotnosti
balení mouky uváděné v kilogramech. Jak se změní variační
koeficient měření, jestliže hmotnost balení uvedeme v gramech?}
{Nezmění se.}{Zvětší se.}{Zmenší se.}{}{}{}}
\LANGsk{
\Question{
Štatistický súbor obsahuje údaje o opakovanom meraní hmotnosti balenia v kilogramoch. Ako sa zmení variačný koeficient merania, ak hmotnosť balenia uvedieme v gramoch?}
{Nezmení sa.}{Zväčší sa.}{Zmenší sa.}{}{}{}}
\LANGes{
\Question{
Statistical file contains repeating measurements of a body mass in kilograms. Find
the change in the coefficient of variation if we convert all data in the file to grams.}
{The coefficient of variation does not change.}{The coefficient of variation becomes bigger.}{The coefficient of variation becomes smaller.}{}{}{}}
\End

\Data\ID{9000153310}
\Updated{2020-07-15 03:07:37}
\Level{B}
\SubArea{Statistics}
\LANGen{
\Question{
Student did repeated measurements of the coefficient of friction,
which is a dimensionless quantity. The mean of the measurements is
\(0.6\) and the coefficient of
variation (relative error) \(10\%\).
From statistics it is known that with a probability nearly
\(100\%\) the
value of the coefficient is in the interval centered at the mean and having radius equal
to the triple of the standard deviation. Find the upper bound of this interval.}
{\(0.78\)}{\(0.18\)}{\(0.42\)}{\(0.66\)}{}{}}
\LANGpl{
\Question{
Student wielokrotnie przeprowadzał pomiary współczynnika tarcia, jest to wartość niewymiarowa. Średnia pomiaru wynosi 
\(0{,}6\), współczynnik zmienności wynosi 
 (błąd względny) \(10\%\).
Ze statystyki wiemy, że jeżeli  prawdopodobieństwo jest bliskie 
\(100\%\) to wartość współczynnika zmienności jest w 
 przedziale  wyśrodkowana w średniej i  posiada promień równy
potrójnej wartości odchylenia standardowego. Wyznacz górną granicę tego przedziału.}
{\(0{,}78\)}{\(0{,}18\)}{\(0{,}42\)}{\(0{,}66\)}{}{}}
\LANGcs{
\Question{
Student měřil koeficient smykového tření (bezrozměrné
číslo). Aritmetický průměr jeho měření byl \(0{,}6\) a relativní chyba
měření (variační koeficient) byla \(10\:\%\). Jaký připouštíme nejvyšší
koeficient tření, jestliže maximální chyba měření (tzv. krajní chyba)
je ve výši trojnásobku směrodatné odchylky měření?}
{\(0{,}78\)}{\(0{,}18\)}{\(0{,}42\)}{\(0{,}66\)}{}{}}
\LANGsk{
\Question{
Študent opakovane meral koeficient šmykového trenia (bezrozmerné číslo). Aritmetický priemer jeho merania bol \(0{,}6\) a variačný koeficient (relatívna chyba merania) bola \(10\%\). Aký najvyšší koeficient trenia pripustíme, ak maximálna chyba merania (tzv. krajná chyba) je vo výške trojnásobku smerodajnej odchýlky merania?}
{\(0{,}78\)}{\(0{,}18\)}{\(0{,}42\)}{\(0{,}66\)}{}{}}
\LANGes{
\Question{
Student did repeated measurements of the coefficient of friction,
which is a dimensionless quantity. The mean of the measurements is
\(0.6\) and the coefficient of
variation (relative error) \(10\%\).
From statistics it is known that with a probability nearly
\(100\%\) the
value of the coefficient is in the interval centered at the mean and having radius equal
to the triple of the standard deviation. Find the upper bound of this interval.}
{\(0.78\)}{\(0.18\)}{\(0.42\)}{\(0.66\)}{}{}}
\End

\Data\ID{1003134409}
\Updated{2020-07-15 03:07:21}
\Level{C}
\SubArea{Statistics}
\LANGen{
\Question{
Twenty-five students of the \( 7 \)th grade took an IQ test and SAT Reasoning Test. Results of the tests are denoted by IQ and SQ in headers of the following cross table. In the table the numbers of students are listed according to the results in both tests, while the results of both tests are classified in intervals. Determine the correlation coefficient between IQ and SQ. Round the result to four decimal places. Use your calculator in Statistics mode to carry out statistical calculations.
\[ \begin{array}{|c|c|c|c|c|} \hline \textbf{SQ \ IQ} & \mathbf{(85;95]} & \mathbf{(95;105]} & \mathbf{(105;115]} & \mathbf{(115;125]} \\\hline \mathbf{(40;60]} & 1 & & & \\\hline \mathbf{(60;80]} & & 10 & 6 & 1 \\\hline \mathbf{(80;100]} & & & 6 & 1 \\\hline \end{array}\]}
{\( 0{.}6086 \)}{\( 0{.}0086 \)}{\( 0{.}9605 \)}{\( -0{.}6806 \)}{}{}}
\LANGpl{
\Question{
Dwudziestu pięciu uczniów klasy \( 7 \) przystąpiło do testu  IQ i  SAT. Wyniki testów znajdują się w tabeli.  W tabeli podano liczbę uczniów według wyników obu testów,  wyniki obu testów zostały uzyskane  w odstępach czasowych.  Określ współczynnik korelacji pomiędzy  IQ i SQ. Zaokrągli wynik do czterech miejsc po przecinku.  Użyj kalkulatora w trybie Statystyki do przeprowadzenia kalkulacji.
\[ \begin{array}{|c|c|c|c|c|} \hline \textbf{SQ \ IQ} & \mathbf{(85;95\rangle} & \mathbf{(95;105\rangle} & \mathbf{(105;115\rangle} & \mathbf{(115;125\rangle} \\\hline \mathbf{(40;60\rangle} & 1 & & & \\\hline \mathbf{(60;80\rangle} & & 10 & 6 & 1 \\\hline \mathbf{(80;100\rangle} & & & 6 & 1 \\\hline \end{array}\]}
{\( 0{,}6086 \)}{\( 0{,}0086 \)}{\( 0{,}9605 \)}{\( -0{,}6806 \)}{}{}}
\LANGcs{
\Question{
Dvacet pět žáků sedmých tříd absolvovalo inteligenční test, jehož výsledkem je tzv. inteligenční kvocient (IQ) a také test všeobecných studijních předpokladů, jehož výsledek označíme SQ. V následující tabulce jsou uvedeny četnosti žáků podle jejich výsledků v obou testech, přičemž výsledky obou testů jsou roztříděné do intervalů. Určete korelační koeficient mezi výsledky inteligenčního testu a testu všeobecných studijních předpokladů. Výsledek zaokrouhlete na čtyři desetinná místa. Pro výpočty použijte správný statistický režim kalkulačky.
\[ \begin{array}{|c|c|c|c|c|} \hline \textbf{SQ \ IQ} & \mathbf{(85;95\rangle} & \mathbf{(95;105\rangle} & \mathbf{(105;115\rangle} & \mathbf{(115;125\rangle} \\\hline \mathbf{(40;60\rangle} & 1 & & & \\\hline \mathbf{(60;80\rangle} & & 10 & 6 & 1 \\\hline \mathbf{(80;100\rangle} & & & 6 & 1 \\\hline \end{array}\]}
{\( 0{,}6086 \)}{\( 0{,}0086 \)}{\( 0{,}9605 \)}{\( -0{,}6806 \)}{}{}}
\LANGsk{
\Question{
Dvadsaťpäť žiakov siedmych tried sa podrobilo inteligenčnému testu, ktorého výsledkom je tzv. inteligenční kvocient (IQ) a aj testu všeobecných študijných predpokladov, ktorého výsledok označíme SQ. V nasledujúcej tabuľke sú zapísané početnosti žiakov podľa ich výsledkov v obidvoch testoch, pričom výsledky obidvoch testov sú  roztriedené do intervalov. Určite korelačný koeficient medzi výsledkami inteligenčného testu a testu všeobecných študijných predpokladov. Výsledok zaokrúhlite na štyri desatinné miesta. Na výpočty použite správny štatistický režim kalkulačky.
\[ \begin{array}{|c|c|c|c|c|} \hline \textbf{SQ \ IQ} & \mathbf{(85;95\rangle} & \mathbf{(95;105\rangle} & \mathbf{(105;115\rangle} & \mathbf{(115;125\rangle} \\\hline \mathbf{(40;60\rangle} & 1 & & & \\\hline \mathbf{(60;80\rangle} & & 10 & 6 & 1 \\\hline \mathbf{(80;100\rangle} & & & 6 & 1 \\\hline \end{array}\]}
{\( 0{,}6086 \)}{\( 0{,}0086 \)}{\( 0{,}9605 \)}{\( -0{,}6806 \)}{}{}}
\LANGes{
\Question{
Twenty-five students of the \( 7 \)th grade took an IQ test and SAT Reasoning Test. Results of the tests are denoted by IQ and SQ in headers of the following cross table. In the table the numbers of students are listed according to the results in both tests, while the results of both tests are classified in intervals. Determine the correlation coefficient between IQ and SQ. Round the result to four decimal places. Use your calculator in Statistics mode to carry out statistical calculations.
\[ \begin{array}{|c|c|c|c|c|} \hline \textbf{SQ \ IQ} & \mathbf{(85;95]} & \mathbf{(95;105]} & \mathbf{(105;115]} & \mathbf{(115;125]} \\\hline \mathbf{(40;60]} & 1 & & & \\\hline \mathbf{(60;80]} & & 10 & 6 & 1 \\\hline \mathbf{(80;100]} & & & 6 & 1 \\\hline \end{array}\]}
{\( 0{.}6086 \)}{\( 0{.}0086 \)}{\( 0{.}9605 \)}{\( -0{.}6806 \)}{}{}}
\End

\Data\ID{1103134408}
\Updated{2021-02-19 01:02:33}
\Level{C}
\SubArea{Statistics}
\IMG{1103134408.pdf}
\LANGen{
\Question{
The values of variables \( x \) and \( y \) are listed in the following table and visualized in the next graph. Calculate the correlation coefficient of \( x \) and \( y \) and round it to four decimal places.
\[ \begin{array}{|c|c|c|c|c|c|} \hline x & 5 & 6 & 7 & 9 & 11 \\\hline y & 3 & 2 &4 & 6 & 8 \\\hline \end{array} \]}
{\( 0{.}9569 \)}{\( 0{.}9659 \)}{\( 0{.}9695 \)}{\( 0{.}9596 \)}{}{}}
\LANGpl{
\Question{
Wartości zmiennych \( x \) i  \( y \) podano w tabeli i na wykresie. Oblicz współczynnik korelacji \( x \) i \( y \) i zaokrągli do czterech miejsc po przecinku.
\[ \begin{array}{|c|c|c|c|c|c|} \hline x & 5 & 6 & 7 & 9 & 11 \\\hline y & 3 & 2 &4 & 6 & 8 \\\hline \end{array} \]}
{\( 0{,}9569 \)}{\( 0{,}9659 \)}{\( 0{,}9695 \)}{\( 0{,}9596 \)}{}{}}
\LANGcs{
\Question{
Vypočítejte koeficient korelace pro znaky \( x \) a \( y \), jejichž hodnoty jsou dány následující tabulkou a zobrazené v grafu. Výsledky zaokrouhlete na čtyři desetinná místa.
\[ \begin{array}{|c|c|c|c|c|c|} \hline x & 5 & 6 & 7 & 9 & 11 \\\hline y & 3 & 2 &4 & 6 & 8 \\\hline \end{array} \]}
{\( 0{,}9569 \)}{\( 0{,}9659 \)}{\( 0{,}9695 \)}{\( 0{,}9596 \)}{}{}}
\LANGsk{
\Question{
Vypočítajte koeficient korelácie pre znaky \( x \) a \( y \), ich hodnoty sú dané nasledujúci tabuľkou a zobrazené v grafe. Výsledky zaokrúhlite na štyri desatinné miesta.
\[ \begin{array}{|c|c|c|c|c|c|} \hline x & 5 & 6 & 7 & 9 & 11 \\\hline y & 3 & 2 &4 & 6 & 8 \\\hline \end{array} \]}
{\( 0{,}9569 \)}{\( 0{,}9659 \)}{\( 0{,}9695 \)}{\( 0{,}9596 \)}{}{}}
\LANGes{
\Question{
Calcula el coeficiente de corelarión entre \( x \) e \( y \), cuyos valores están en la tabla y dibujadas en la gráfica.Aproxima los resultados a cuatro cifras decimales.
\[ \begin{array}{|c|c|c|c|c|c|} \hline x & 5 & 6 & 7 & 9 & 11 \\\hline y & 3 & 2 &4 & 6 & 8 \\\hline \end{array} \]}
{\( 0{,}9569 \)}{\( 0{,}9659 \)}{\( 0{,}9695 \)}{\( 0{,}9596 \)}{}{}}
\End

\Data\ID{1103134410}
\Updated{2020-07-15 03:07:21}
\Level{C}
\SubArea{Statistics}
\IMG{1103134410.pdf}
\LANGen{
\Question{
The heights (\(\mathrm{cm}\)) of ten boys and their best performances (\(\mathrm{cm}\)) in standing long jump at the international athletic championship are listed in the following table. Determine the correlation coefficient \( r \) between height of boys and their longest jump. You can use a statistics mode of your calculator to carry out statistical computations. Round the result to four decimal places. 
Based on the following scatterplot and the correlation coefficient interpret the strength of linear relationship between analysed variables.
\[ \begin{array}{|c|c|c|c|c|c|} \hline 
\textbf{Boy's height (cm)} & 189 & 175 & 187 & 183 & 174 \\\hline \textbf{Length of the jump (cm)} & 231 & 207 & 214 & 223 & 202  \\\hline \\\hline
\textbf{Boy's height (cm)} & 193 & 179 & 169 & 186 & 183 \\\hline \textbf{Length of the jump (cm)} & 242 & 229 & 190 & 226 & 212 \\\hline \end{array} \]}
{strong linear relationship: \( r = 0{.}8628 \)}{moderate linear relationship: \( r = 0{.}5542 \)}{moderate linear relationship: \( r = 0{.}7444 \)}{strong linear relationship: \( r = 0{.}9289 \)}{}{}}
\LANGpl{
\Question{
Wzrost (angl. Height) chłopców oraz  ich najlepsze skoki w dal (angl. Length of the jump) na międzynarodowych zawodach zostały podane w tabeli.  Określ współczynnik korelacji \( r \) pomiędzy wysokością chłopców a ich najdłuższym skokiem. Możesz użyć trybu Statystyka w  kalkulatorze, aby przeprowadzić obliczenia statystyczne. Zaokrąglij wynik do czterech miejsc po przecinku.
Na podstawie poniższego wykresu rozrzutu i współczynnika korelacji zinterpretuj  siłę  zależności liniowej  pomiędzy analizowanymi zmiennymi.

\[ \begin{array}{|c|c|c|c|c|c|} \hline 
\textbf{Wzrost (cm)} & 189 & 175 & 187 & 183 & 174  \\\hline \textbf{Długość skoku (cm)} & 231 & 207 & 214 & 223 & 202  \\\hline \\\hline
\textbf{Wzrost (cm)} & 193 & 179 & 169 & 186 & 183 \\\hline \textbf{Długość skoku (cm)} & 242 & 229 & 190 & 226 & 212 \\\hline \end{array} \]}
{silna zależność liniowa: \( r = 0{,}8628 \)}{średnia zależność liniowa: \( r = 0{,}5542 \)}{średnia zależność liniowa: \( r = 0{,}7444 \)}{silna zależność liniowa: \( r = 0{,}9289 \)}{}{}}
\LANGcs{
\Question{
V tabulce jsou uvedeny výšky deseti chlapců (angl. Height) a jejich nejlepší výkony ve skoku z místa do dálky (angl. Length of the jump) na mezinárodních závodech. Určete korelační koeficient \( r \) mezi výškou skokana a jeho výkonností v této disciplíně. Výsledek zaokrouhlete na čtyři desetinná místa. Na základě bodového grafu na obrázku a hodnoty korelačního koeficientu posuďte míru lineární závislosti mezi výškou skokana a délkou jeho skoku. Pro výpočty použijte správný statistický režim kalkulačky.
\[ \begin{array}{|c|c|c|c|c|c|} \hline 
\textbf{Výška žáka (cm)} & 189 & 175 & 187 & 183 & 174 \\\hline \textbf{Délka skoku (cm)} & 231 & 207 & 214 & 223 & 202  \\\hline \\\hline
\textbf{Výška žáka (cm)} & 193 & 179 & 169 & 186 & 183 \\\hline \textbf{Délka skoku (cm)} & 242 & 229 & 190 & 226 & 212 \\\hline 
\end{array} \]}
{silná lineární závislost: \( r = 0{,}8628 \)}{středně silná lineární závislost: \( r = 0{,}5542 \)}{středně silná lineární závislost: \( r = 0{,}7444 \)}{silná lineární závislost: \( r = 0{,}9289 \)}{}{}}
\LANGsk{
\Question{
V tabuľke sú zaznamenané výšky (angl. Height) desiatich chlapcov a ich najlepšie výkony v skoku z miesta do diaľky (angl. Length of the jump) na medzinárodných pretekoch. Určte korelačný koeficient \( r \) medzi výškou skokana a jeho výkonnosťou v tejto disciplíne. Výsledok zaokrúhlite na štyri desatinné miesta. Na základe bodového grafu na obrázku a hodnoty korelačného koeficientu posúďte mieru lineárnej závislosti medzi výškou skokana a dĺžkou jeho skoku. Na výpočty použite správny štatistický režim kalkulačky.
\[ \begin{array}{|c|c|c|c|c|c|} \hline 
\textbf{Výška žiaka (cm)} & 189 & 175 & 187 & 183 & 174  \\\hline \textbf{Dĺžka skoku (cm)} & 231 & 207 & 214 & 223 & 202  \\\hline \\\hline
\textbf{Výška žiaka (cm)} & 193 & 179 & 169 & 186 & 183 \\\hline \textbf{Dĺžka skoku (cm)} & 242 & 229 & 190 & 226 & 212 \\\hline \end{array} \]}
{silná lineárna závislosť: \( r = 0{,}8628 \)}{stredne silná lineárna závislosť: \( r = 0{,}5542 \)}{stredne silná lineárna závislosť: \( r = 0{,}7444 \)}{silná lineárna závislosť: \( r = 0{,}9289 \)}{}{}}
\LANGes{
\Question{
The heights (\(\mathrm{cm}\)) of ten boys and their best performances (\(\mathrm{cm}\)) in standing long jump at the international athletic championship are listed in the following table. Determine the correlation coefficient \( r \) between height of boys and their longest jump. You can use a statistics mode of your calculator to carry out statistical computations. Round the result to four decimal places. 
Based on the following scatterplot and the correlation coefficient interpret the strength of linear relationship between analysed variables.
\[ \begin{array}{|c|c|c|c|c|c|} \hline 
\textbf{Boy's height (cm)} & 189 & 175 & 187 & 183 & 174 \\\hline \textbf{Length of the jump (cm)} & 231 & 207 & 214 & 223 & 202  \\\hline \\\hline
\textbf{Boy's height (cm)} & 193 & 179 & 169 & 186 & 183 \\\hline \textbf{Length of the jump (cm)} & 242 & 229 & 190 & 226 & 212 \\\hline \end{array} \]}
{strong linear relationship: \( r = 0{.}8628 \)}{moderate linear relationship: \( r = 0{.}5542 \)}{moderate linear relationship: \( r = 0{.}7444 \)}{strong linear relationship: \( r = 0{.}9289 \)}{}{}}
\End

\Data\ID{1103134411}
\Updated{2020-07-15 03:07:10}
\Level{C}
\SubArea{Statistics}
\LANGen{
\Question{
The following scatter plots represent visualisations of relationships of two continuous variables. Choose the plot, which represents visualisation of a relationship of two variables with the highest absolute value of the correlation coefficient.}
{}{}{}{}{}{}}
\LANGpl{
\Question{
Poniższe wykresy rozproszenia reprezentują wizualizacje zależności dwóch zmiennych ciągłych. Wybierz wykres, który reprezentuje wizualizację relacji dwóch zmiennych o najwyższej wartości bezwzględnej współczynnika korelacji.}
{}{}{}{}{}{}}
\LANGcs{
\Question{
Na obrázcích jsou zobrazené bodové grafy (korelační pole), představující závislost dvou znaků. Vyberte z těchto grafů ten, na kterém je zobrazena závislost dvou znaků s největší absolutní hodnotou korelačního koeficientu.}
{}{}{}{}{}{}}
\LANGsk{
\Question{
Na obrázkoch sú zobrazené bodové grafy (korelačné pole), predstavujúce závislosť dvoch znakov. Vyberte z týchto grafov ten, na ktorom je zobrazená závislosť dvoch znakov s najväčšou absolútnou hodnotou korelačného koeficientu.}
{}{}{}{}{}{}}
\LANGes{
\Question{
The following scatter plots represent visualisations of relationships of two continuous variables. Choose the plot, which represents visualisation of a relationship of two variables with the highest absolute value of the correlation coefficient.}
{}{}{}{}{}{}}

\IMGanswer{1103134411a_0.pdf}
\IMGanswer{1103134411b_0.pdf}
\IMGanswer{1103134411c_0.pdf}
\IMGanswer{1103134411d_0.pdf}\End

\Data\ID{9000153302}
\Updated{2020-07-15 03:07:37}
\Level{C}
\SubArea{Statistics}
\LANGen{
\Question{
A student performed repeated measurements of a length (in meters) and evaluated
the main statistical characteristics: mean, standard deviation, variance and the
coefficient of variation. Which of these characteristics is dimensionless?}
{coefficient of variation}{variance}{standard deviation}{mean}{}{}}
\LANGpl{
\Question{
Student przeprowadził wielokrotne pomiary długości (w metrach) i ocenił
główne cechy statystyczne: średnią, odchylenie standardowe, wariancję i
współczynnik zmienności. Która z tych cech nie posiada jednostki miary?}
{współczynnik zmienności}{wariancja}{odchylenie standardowe}{średnia}{}{}}
\LANGcs{
\Question{
Student opakovaně měřil délku tělesa (v metrech). Která z níže
uvedených charakteristik měření je bez jednotky (bezrozměrné číslo)?}
{variační koeficient}{rozptyl}{směrodatná odchylka}{aritmetický průměr}{}{}}
\LANGsk{
\Question{
Študent opakovane meral dĺžku telesa (v metroch). Ktorá z nižšie uvedených charakteristík merania je bez jednotky (bezrozmerné číslo)?}
{variačný koeficient}{rozptyl}{smerodajná odchýlka}{aritmetický priemer}{}{}}
\LANGes{
\Question{
A student performed repeated measurements of a length (in meters) and evaluated
the main statistical characteristics: mean, standard deviation, variance and the
coefficient of variation. Which of these characteristics is dimensionless?}
{coefficient of variation}{variance}{standard deviation}{mean}{}{}}
\End

\Data\ID{9000153303}
\Updated{2020-07-15 03:07:37}
\Level{C}
\SubArea{Statistics}
\LANGen{
\Question{
A student performed repeated measurements of a length (in meters) and evaluated
the main statistical characteristics: mean, standard deviation, variance and
the coefficient of variation. Unit of which of these characteristics is meter?}
{the mean and the standard deviation}{only the variance}{only the standard deviation}{only the mean}{the standard deviation and the variance}{the standard deviation, the variance and the coefficient of variation}}
\LANGpl{
\Question{
Student przeprowadził wielokrotne pomiary długości (w metrach) i ocenił
główne cechy statystyczne: średnią, odchylenie standardowe, wariancję i
współczynnik zmienności. Wskaż cechę statystyczną, której jednostką miary jest metr?}
{średnia i  odchylenie standardowe}{tylko wariancja}{tylko odchylenie  standardowe}{tylko średnia}{odchylenie  standardowe i wariancja}{odchylenie  standardowe,  wariancja i współczynnik zmienności}}
\LANGcs{
\Question{
Student opakovaně měřil délku tělesa (v metrech). Které ze standardně
uváděných charakteristik měření (aritmetický průměr, směrodatná
odchylka, rozptyl, variační koeficient) mají jednotku metr?}
{aritmetický průměr a směrodatná odchylka}{pouze rozptyl}{pouze směrodatná odchylka}{pouze aritmetický průměr}{směrodatná odchylka a rozptyl}{směrodatná odchylka, rozptyl a variační koeficient}}
\LANGsk{
\Question{
Študent opakovane meral dĺžku telesa (v metroch). Ktoré zo štandardne uvádzaných charakteristík (aritmetický priemer, smerodajná odchýlka, rozptyl, variačný koeficient) majú jednotku meter?}
{aritmetický priemer a smerodajná odchýlka}{len rozptyl}{len smerodajná odchýlka}{len aritmetický priemer}{smerodajná odchýlka a rozptyl}{smerodajná odchýlka, rozptyl a variačný koeficient}}
\LANGes{
\Question{
A student performed repeated measurements of a length (in meters) and evaluated
the main statistical characteristics: mean, standard deviation, variance and
the coefficient of variation. Unit of which of these characteristics is meter?}
{the mean and the standard deviation}{only the variance}{only the standard deviation}{only the mean}{the standard deviation and the variance}{the standard deviation, the variance and the coefficient of variation}}
\End
